% Generated by scripts/build_guide.py from reviewed-section paragraphs.
\def\EditionMode{3}
\def\EditionLabel{A guide for the reader}
\def\EditionStatus{WORKING GUIDE: 58 completed section/front-matter paragraphs; 233 of 706 source pages reviewed}
\documentclass[10pt,oneside]{book}
\usepackage{fontspec}
\setmainfont{Latin Modern Roman}
\setsansfont{Latin Modern Sans}
\usepackage{amsmath,amssymb,mathtools}
\usepackage{unicode-math}
\setmathfont{Latin Modern Math}
\usepackage{xcolor,tikz-cd,lua-ul,marginnote,pdflscape}
\usetikzlibrary{arrows.meta,decorations.pathmorphing,calc,shapes.geometric}
\tikzcdset{ampersand replacement=\&}
\tikzcdset{every label/.append style={text=black}}
\usepackage[a4paper,left=18mm,right=45mm,top=18mm,bottom=22mm,marginparwidth=28mm,marginparsep=7mm]{geometry}
\usepackage[unicode,colorlinks=true,urlcolor=black,linkcolor=black]{hyperref}
\usepackage{bookmark}
\setlength{\parindent}{0pt}
\setlength{\parskip}{4pt plus 1pt minus 1pt}
\setlength{\emergencystretch}{1.5em}
\raggedbottom
\definecolor{PSHand}{RGB}{25,58,132}
\definecolor{PSDelete}{RGB}{146,37,37}
\definecolor{PSTypedOld}{RGB}{95,95,95}
\ifnum\EditionMode=2\colorlet{blue}{black}\else\colorlet{blue}{PSHand}\fi
\RenewDocumentCommand{\textcircled}{m}{\tikz[baseline=(PSCircle.base)]\node[draw,circle,inner sep=.5pt,line width=.4pt](PSCircle){#1};}
\newif\ifPSInApparatus
\newcounter{psnote}
\newcounter{psvariant}
\newcounter{pslinkmark}
\newif\ifPSInLink
\newsavebox{\PSNoteBox}
\newdimen\PSNoteLimit
{\footnotesize\global\PSNoteLimit=3\baselineskip}
\newcommand{\CurrentScanURL}{}
\newcommand{\CurrentKey}{}
\newcommand{\CurrentScan}{}
\newcommand{\pssource}[3]{\psflushmargins\clearpage\gdef\CurrentKey{#1}\gdef\CurrentScan{#2}\gdef\CurrentScanURL{#3}\setcounter{psvariant}{0}\typeout{PS-SOURCE: #1}}
\newcommand{\pslinkmarker}[2]{\stepcounter{pslinkmark}\latelua{texio.write_nl('log','PS-LINK-SHIP|#1|\CurrentKey|\arabic{pslinkmark}|#2')}}
\newcommand{\pslink}[2]{\ifPSInLink\pslinkmarker{#1}{nested}#2\else\begingroup\PSInLinktrue\href{\CurrentScanURL}{\pslinkmarker{#1}{outer}#2}\endgroup\fi}
\newcommand{\scanpage}[1]{\ifnum\EditionMode=2\marginnote{\scriptsize\sffamily\pslink{scanpage}{scan #1}}\else\par\noindent{\small\sffamily\pslink{scanpage}{[scan #1]}}\par\medskip\fi}
\newcommand{\mspage}[1]{\ifnum\EditionMode=2\marginnote{\scriptsize\sffamily p.\,#1}\else{\small\sffamily[p.\,#1]}\fi}
\newcommand{\hand}[1]{{\ifnum\EditionMode=2\color{black}\else\color{PSHand}\fi #1}}
\let\over\relax
\NewDocumentCommand{\over}{O{} m m}{{\textcolor{PSTypedOld}{#2}\ensuremath{^{\!\nearrow}}\if\relax\detokenize{#1}\relax\hand{#3}\else #3\fi}}
\NewDocumentCommand{\inserted}{O{} m}{{\if\relax\detokenize{#1}\relax\hand{\ensuremath{\langle}#2\ensuremath{\rangle}}\else\ensuremath{\langle}#2\ensuremath{\rangle}\fi}}
\newcommand{\deleted}[1]{{\color{PSDelete}\ifmmode\text{[del.: }#1\text{]}\else[\textsf{del.:} #1]\fi}}
\newcommand{\xout}[1]{{\color{PSTypedOld}\ifmmode\text{[x: }#1\text{]}\else[\textsf{x:} #1]\fi}}
\let\PSMathUnderline\underline
\RenewDocumentCommand{\underline}{+m}{\ifmmode\PSMathUnderline{#1}\else\underLine{#1}\fi}
\NewDocumentCommand{\uline}{+m}{\underline{#1}}
\NewDocumentCommand{\uuline}{+m}{\ifmmode\PSMathUnderline{\PSMathUnderline{#1}}\else\underLine{\underLine{#1}}\fi}
\newunderlinetype\PSWaveLine{\leaders\hbox{\lower.7ex\hbox{\tikz[x=1pt,y=1pt,baseline=0pt]{\draw[line width=.3pt] (0,0) .. controls (1,1) and (2,1) .. (3,0) .. controls (4,-1) and (5,-1) .. (6,0);}}}}
\NewDocumentCommand{\uwave}{+m}{{\PSWaveLine #1}}
\newcommand{\unsure}[1]{\pslink{unsure}{[#1\ifmmode\,?\else\,?\fi]}}
\newcommand{\illegible}[1]{\ifmmode\text{\pslink{illegible}{[illeg.\ #1]}}\else\pslink{illegible}{[illeg.\ #1]}\fi}
\NewDocumentCommand{\diagram}{+m}{\ifmmode\vcenter{\hbox{#1}}\else\par\medskip{\centering #1\par}\medskip\psflushmargins\fi}
\newenvironment{diagramblock}{\par\begin{center}}{\end{center}\par}
\newcommand{\PSMathNotes}{}
\makeatletter
\NewDocumentCommand{\psplacednote}{m +m}{%
      \sbox{\PSNoteBox}{\parbox{\marginparwidth}{\footnotesize\raggedright\hand{\textsf{m#1.} #2}}}%
      \ifdim\dimexpr\ht\PSNoteBox+\dp\PSNoteBox\relax<\PSNoteLimit
        \leavevmode\marginpar{\usebox{\PSNoteBox}}%
      \else
        \par\nobreak\smallskip\hrule\nobreak\smallskip
        {\footnotesize\textsf{Marginal note m#1.}\quad\hand{#2}\par}
        \nobreak\smallskip\hrule\smallskip
      \fi}
\newcommand{\psflushmargins}{\ifmmode\else\ifinner\else\ifx\PSMathNotes\@empty\else\let\PSTempNotes\PSMathNotes\global\let\PSMathNotes\@empty\PSTempNotes\fi\fi\fi}
\NewDocumentCommand{\margin}{+m}{%
  \ifPSInApparatus\ifmmode\hand{\text{[marg.: }#1\text{]}}\else\hand{\textsf{[marg.: }#1\textsf{]}}\fi
  \else\stepcounter{psnote}%
    \ifmmode
      \text{\textsuperscript{m\arabic{psnote}}}%
      \begingroup\edef\PSQueue{\endgroup\noexpand\g@addto@macro\noexpand\PSMathNotes{\noexpand\psplacednote{\arabic{psnote}}{\unexpanded{$#1$}}}}\PSQueue
    \else
      \textsuperscript{m\arabic{psnote}}\psplacednote{\arabic{psnote}}{#1}%
    \fi
  \fi}
\AtBeginDocument{\everymath\expandafter{\the\everymath\aftergroup\psflushmargins}\everydisplay\expandafter{\the\everydisplay\aftergroup\psflushmargins}}
\makeatother
\newcommand{\psnumber}[2]{\expandafter\gdef\csname psnum@#1\endcsname{#2}}
\NewDocumentCommand{\var}{+m +m m m}{\ifnum\EditionMode=0 #1\else\ifnum\EditionMode=1\ifmmode\text{\textsuperscript{\tiny\csname psnum@#4\endcsname}}\else\textsuperscript{\tiny\csname psnum@#4\endcsname}\fi\fi #2\fi}
\newcommand{\psapparatusstart}{\par\bigskip\hrule\smallskip\PSInApparatustrue\begingroup\footnotesize\textsf{Apparatus for scan \CurrentScan}\par}
\newcommand{\psapparatusend}{\par\endgroup\PSInApparatusfalse}
\NewDocumentCommand{\psentry}{m m m +m +m m}{\par\noindent\textbf{#2. #3.}\quad\textit{Typescript:} #4\quad\textit{Corrected:} \if\relax\detokenize{#5}\relax\textnormal{[omitted]}\else #5\fi\quad #6\par}
\hyphenation{Gro-then-dieck abel-ian-iza-tion}

\begin{document}
\hypersetup{pageanchor=false}
\begin{titlepage}\raggedright
{\Large Alexander Grothendieck\par}\vspace{12mm}
{\Huge Pursuing Stacks\par}\medskip{\LARGE 1983\par}\vspace{12mm}
{\Large A guide for the reader\par}\vspace{8mm}
{\large\EditionStatus\par}\vfill
Université de Montpellier, Grothendieck archive, cotes 134-2--134-8.\par
Including the appendix translation by R. Brown, cote 134-2, sites 115--143.\par
\end{titlepage}\hypersetup{pageanchor=true}
\chapter*{Editorial statement}

AI-transcribed by Claude Opus 5.5 (Anthropic, effort high); checked against the scans, and this edition built, by GPT-6 Astra (OpenAI, ChatGPT Pro mode); prompts and records by Claude (Anthropic).

The witness consulted and not reproduced is the LaTeX edition of Mateo Carmona with the collaboration of Ulrik Buchholtz (arXiv 2111.01000 v2, 2021), extended from the 2015 edition by ``the scrivener'', CC0 1.0.

The translation in the appendix is R. Brown's, prepared with assistance from L. Breen and J-L. Loday, as his note says; the rights in it belong to its translator or his heirs.

The rights in the original work belong to its author or the author's representatives or heirs. Where we hold rights, in the AI-produced transcription, translation, corrections and typesetting, they are dedicated to the public domain under CC0 1.0. This edition is published for educational, non-profit purposes, and we claim no rights we do not hold. If you hold rights in this work and object to its publication here, please open a GitHub issue and we will take it down.


\section*{Using this guide}
Each paragraph describes a section already checked against both supplied scan halves. It identifies the section's task, terminology, dependencies, revisions and consequential witness differences. Mathematical cautions are editorial observations, distinguished from the author's statements; they do not alter the readers. Archive sites and internal typescript folios are kept distinct, especially where the initial letter has been reordered or renumbered. The account is provisional wherever the relevant later section has not yet been reviewed.

The guide is incomplete while source collation continues. No paragraph for an unreviewed section has been supplied from the arXiv witness alone. The completed paragraphs follow the original section order; gaps in the current working guide are intentional and will close as their sections are checked.
\chapter*{Section guide}
\subsection*{Front matter: plans, contents, and dedication}
The front matter (cote 134-2, sites 1--13) should be read as evidence of a developing book: the title page and chapter-dependency diagram precede several layers of contents, including cancelled headings, renumbered sections, and an entirely handwritten replacement page. The diagram uses the author's chapter labels and places derivators within the projected sequence; its heavily retraced connections should not be treated as a settled logical dependency chart. The cancelled chapter-IV block on site 7 and the replacement handwritten contents on site 8 show the reorganization of the asphericity discussion. Sites 9--12 then project the progression through homology and abelianization, schematization, and linearization: the proposed comparison between schematic and discrete homotopy theory leads to the ``breakdown of an idyllic picture'' of section 119, then to revised programs and questions about the six operations in sections 124 and 134--139. These titles announce later developments; they do not themselves establish the claims. In particular, the pair of arrows listed for section 121 points from homotopy types to abelian homotopy types above and back below, with handwritten arrows and accents over typed objects and letters. Several continuations on site 11---\emph{topos}, \emph{bundles}, and \emph{valences}---are typing guided by handwritten curves, a material distinction lost in the initial transcription. The doubled \emph{ff} in the title for section 131 is visibly typed. The French dedication on site 13 retains a cancelled address to former pupils above the dedication to Dodek Shapiro, together with its cancelled and unresolved words. It remains French in all readers. The contents' cancellations and the final handwritten ``to be completed'' are particularly useful reminders that the source records successive intentions, including abandoned formulations, rather than a retrospectively regularized exposition (Grothendieck, 1983, cote 134-2, sites 1--13).

\subsection*{1. Foundational questions and tame topology (letter pp. 1--1')}
The letter dated 19 February 1983 starts from the papers sent by Ronnie Brown and recalls the 1975 correspondence with Larry Breen as a program joining homotopical and homological algebra over topoi. Before pursuing that program, Grothendieck argues for an even more urgent foundation for ``tame'' topology and stratified tame spaces, contrasting the conceptual difficulty with the inertia of established habits. The next section returns to the homotopical program; section 14 later resumes the criticism of the category of topological spaces and its inadequacy for isotopy and stratification. The opening is therefore a statement of priorities, not a theorem about a proposed category. Its cancellation history matters: on cote 134-2, site 14, the addressee is crossed out, and the long passage about the recipient's axiomatic homotopical algebra and lecture notes is cancelled and replaced by a shorter handwritten phrase. The arXiv witness prints that cancelled continuation as live text. It also substitutes ``as for schemes'' for the clearly handwritten ``or for schemes''. The section's descriptive title in the witness is absent from the scanned opening, which gives only its number. The continuation on site 15 concludes the comparison with programs of homological, homotopical, or topological algebra.

\subsection*{2. Higher groupoids and coherent composition (letter, sites 15--18)}
Returning to the Breen program recalled in section1, Grothendieck proposes that $n$-truncated homotopy types be described by $n$-groupoids, with a fundamental $n$-groupoid, geometric realization, and an $n$-equivalence comparison; relativization over a topos should yield $n$-gerbs and $n$-stacks as noncommutative cohomological coefficients. The assertion that every $n$-groupoid should be $n$-equivalent to a fundamental one is present on site15 but omitted from the arXiv witness, which joins the surrounding clauses. The attempted globular description then introduces sets $F_i$, source, target, identity and inversion maps, and composition at different levels. Strict associativity is rejected as insufficient for the intended homotopy types; equalities are to become specified homotopies, whose own coherence leads from ``secondary'' to ``ternary'' and further structure, illustrated by the pentagon at the start of site18. Section3 asks whether this can support a model-category description; sections18--19 later revise the coherator and realization proposals. The internal reading order here is site15, then site17 (centered folio3), then site16 (folio4), then site18 (folio5); the readers preserve the archive's different order. Several formulas require caution. On site17 the first factor of the composition domain has an overprinted $t$, with the underlying letter still doubtful, where the witness gives $s$; the marginal identity notation is $id_a$, not $d_a$. The handwritten target formula ends with $s^i_1(u)$ in its second right-hand factor, an apparent source/target slip. On site16 the inversion formula places separate inversion marks on both left-hand factors and appears tautological, although the first right-hand letter remains overwritten and doubtful; the witness substitutes a question mark for that right side. The handwritten interchange range also permits equal levels, $1\leq\alpha\leq\nu$: for a one-object group at equal levels, the displayed interchange would force the two middle factors to commute, so this range is too broad for arbitrary fundamental groups. These source readings remain unchanged. The author replaces ``implying'' by ``expressing'' beneath the compatibility conditions, a correction the witness does not apply. The resulting discussion is a developing proposal for coherence, and these local defects should not be treated as established definitions or silently repaired.

\subsection*{3. The model-category question (letter, sites 18--19)}
Grothendieck asks Quillen whether an appropriate notion of higher category already exists, recalling a description by multiple semisimplicial sets but distinguishing strict associativity from the coherent structure sought in section 2. He proposes two connections with homotopical algebra: the conjectural category of infinity-groupoids should model ordinary homotopy types, and a fundamental infinity-groupoid construction should be available in any model category from a cylinder object of its final object. The latter is presented as informative chiefly in the nonpointed case, and the handwritten ending explicitly identifies the cylinder as playing the role of the unit segment. Section 4 immediately reopens the skeletal data of that proposed construction; section 14 returns to the limitations of topological-space models, and the later interval discussion develops the cylinder intuition. These are proposals and questions, without a completed definition or asserted model-structure theorem. On site 18, the handwritten replacement of ``your'' by ``Quillen's'' in the reference to axiomatics is not applied by the arXiv witness, and its ``contents'' misreads the typed ``comments''. Site 19 closes the first letter with a handwritten signature and is numbered 5-prime on the scan, where the witness labels it 3-prime. The ensuing dated continuation on site 20 begins a substantive revision rather than simply continuing an established definition.

\subsection*{4. Skeletal data and successive structure (letter, sites 20--21)}
In the continuation dated 20 February, Grothendieck revises the undifferentiated ``primary'' structure proposed in section 2: the skeletal datum is only the family of sets with source and target maps, while inversion, degeneration and first composition belong to structure added subsequently. Equivalences of ordinary groupoids motivate replacing strict involutions by pseudoinvolutions, and choices of interval maps explain why degeneracies and path composition should also be nonprimitive. Godement's next composition is then secondary, and the further composition levels successively ternary and higher. This ordering prepares the tower of universal categories in section 5 and makes explicit a change of mind: the text says that the preceding day's operations were hastily put into one pot. Two details deserve care. On site 20 the formal maps $s_o,t_o$ are described as coming from $F_0$ to $F_{-1}$ but are assigned ``$i=1$''; the evident index mismatch is retained. On site 21 the handwritten relocation yields ``are rather'', while the witness retains ``rather are''. The replacement ``the real'' on that page is typed above the line, not handwritten as the initial transcription classified it. The source says ``all these operations'', where the witness shortens this to ``all the operations'', and retains the two handwritten marginal composition symbols absent from that witness. The overwritten folio at the head of site 21 remains uncertain between 7 and 4; the archive order and all scan links are unaffected.

\subsection*{5. A tower of universal categories (letter, sites 21--24)}
Grothendieck proposes a sequence of categories $C_0,C_1,\ldots$ with common objects and progressively added arrows, whose direct limit should classify the proposed infinity-groupoid structures by functors preserving specified limits. The discussion first compares universal structures in categories with finite products, finite inverse limits, and the stronger exactness available in topoi; it then discards the separate finite-product requirement while writing site 23, restricting the common objects to ``standard'' iterated fiber products determined by source and target. This is the categorical organizing device promised in section 4. Section 6 gives it a topological model; the later coherator discussion revises how the successive operations should be adjoined. The text expressly presents the strict increase at every stage as a conviction still needing proof. On site 23 the letter in ``exist in F'' is visibly F, although the ambient category is C and the witness prints C; the source letter is retained as an apparent local slip. The iterated-source indices on that page contain typed 1-shaped bases with handwritten ell/chi-like corrections and remain marked, rather than being settled solely by the witness's ell. The live insertion says that the word ``iterated'' will henceforth be dropped, while the cancelled marginal version says ``qualification''; the witness combines the latter wording with the live text. On site 24 the marginal note assigns order ell to the composition bearing lower index ell, where the witness gives that composition lower index 2. Some of this note's lettering and cancellation status remain doubtful, but the difference of indices should not be silently harmonized. The section's revisions distinguish a proposed universal construction from a finished coherence theorem.

\subsection*{6. Discs and a proposed universal co-groupoid (letter, sites 24--25)}
The tower of section 5 receives its motivating topological interpretation: a space $X$ supplies sets of maps from discs $D_i$, and the source and target maps are induced contravariantly by the two boundary embeddings. Compositions and further operations are to arise from maps between discs and their prescribed amalgamated sums, producing a co-infinity-groupoid whose transform under $h_X$ is the desired infinity-groupoid. Grothendieck then proposes that the category $B_\infty$ generated by these cells, sums and arrows should be dual to the universal category $C_\infty$. The expected universality would make the category independent of the successive arbitrary choices up to unique isomorphism, while he insists that the formal construction of $C_n$ remains meaningful without that expectation. Section 7 begins to specify the disc maps and their coglobular identities. Site 25 records a significant change of mind: the confident parenthesis ``I believe it is'' acquires handwritten exclamation marks, and a marginal addition dated 23 February says ``But no longer now'' and refers to page 10, section 11. The crossing pencil hatching leaves the note's cancellation status uncertain; the readers retain that doubt. The arXiv witness paraphrases this retraction, omits its explicit page reference, and prints the cancelled article before ``various'' amalgamated cells. On site 24 the unit-segment formula retains the typed interlinear $I=$ and the live $D_1$, an identification shortened in the witness. These distinctions preserve both the proposed model and the author's later withdrawal of confidence in its universality.

\subsection*{14. Homotopical algebra and the limits of present topology (pp. 1--2)}
Grothendieck opens the continuation of the letter to Quillen by distinguishing the study of spaces of maps and their successive homotopies from the study of spaces of homeomorphisms, immersions, embeddings and fibrations. His terms ``wildness'', ``isotopy'', ``stratified structures'' and ``tubular neighbourhoods'' delimit problems for which the familiar category of topological spaces seems poorly adapted; the contrast with piecewise linear structures is that their maps are too sparse, notably because real-valued piecewise linear functions are not closed under multiplication. This develops the complaint recalled explicitly from the letter, p.~1, and supplies the contrast on which section~15 builds its more favorable account of localization and homotopy types. The interlinear revisions narrow the opening description to ``maps between spaces'' and replace ``numerical'' by $\mathbb R$-valued (p.~2). The witness differs in two mathematically relevant places: the typescript says ``pro-topoi'' on p.~1, where arXiv has ``pretopoi'', and arXiv retains ``numerical'' without applying its handwritten replacement on p.~2. These readings are established by the scan; the edition does not identify the two notions or silently repair either authorial formulation.

\subsection*{15. Models of homotopy types and the place of stacks (pp. 2--4)}
The argument passes from the failure of the Hauptvermutung as a literal assertion about triangulations to its much weaker homotopical analogue: appropriate localizations of topological spaces and semi-simplicial sets yield the same homotopy category. Grothendieck compares CW-spaces, Kan complexes and Serre's cubical complexes, then proposes stacks, previously called ``$\infty$-groupoïdes'', as a third description particularly close to the homotopy groups themselves. His analogy with chain complexes and the Dold--Puppe equivalence makes the intended role of noncommutative structure explicit; section~16 immediately generalizes the motivation from spaces to topoi. The distinction between describing and computing invariants matters: on p.~3 he twice strikes ``description'' and writes ``computation'', corrections which the arXiv witness does not apply. On p.~4 a marginal note dated 3.5 begins ``This is an error'' and retracts the claim that stack cohomology has no immediate description from the primitive structure, referring instead to the nerve and to forgetting structure. Broad pencil strokes across this note make its final cancellation status uncertain, and one verb remains doubtful; both uncertainties are preserved. The arXiv footnote is substantially longer than this marginal wording and must not be substituted for it as though it were a literal transcription.

\subsection*{16. Stacks over topoi and noncommutative cohomological algebra (pp. 4--6)}
The motivation for stacks is enlarged from section~15's models of homotopy types to coefficients for cohomology over an arbitrary topos. Starting with Giraud's $0$-, $1$- and $2$-dimensional theory, Grothendieck treats $n$-stacks as coefficients whose ``integration'' is evaluation at the final object, with the homotopy groups of the resulting stack supplying cohomology after a dimension shift. His letters to Breen and Mme Raynaud's work on Lefschetz comparisons provide the geometric motivation: higher noncommutative cohomology should unify homological and homotopical comparison statements. The passage to infinite-order stacks is presented as a qualitative enlargement of perspective, and the following section asks how much of the homotopy-theoretic side ordinary categories can already model. The scans matter for the formulas: p.~4 has $H^i$, not the witness's $K^i$; p.~5 has $R\Gamma F_*$ and complexes $F_*$, where the supplied arXiv comparison reports a missing $F$ and a $j$ in place of the star. The handwritten qualification ``[non commutative]'' on p.~5 must also remain visible. These are witness differences, not authorial mathematical corrections introduced by this edition.

\subsection*{16$'$. Categories and the ``impressive bunch'' of models (pp. 7--9)}
Grothendieck associates a homotopy type to a small category both through its presheaf topos and through its nerve, comparing the first construction with Verdier and Artin--Mazur pro-homotopy types. He recalls and rejects his former guess that passage to the enveloping groupoid preserves this type: Quillen had pointed out that ordinary categories realize arbitrary homotopy types. This corrects the initial groupoid intuition while continuing sections~15--16's search for models, and motivates the later study of the ``impressive bunch'' and of coherators. A second proposed shortcut is itself explicitly retracted: the margin of p.~7, dated 5.3, says that the asserted weak equivalence $K\longrightarrow NN'(K)$ is false and directs the reader to section~24, pp.~21--22. The note is shorter than its arXiv counterpart; its call number is (18), and the section has been renumbered $16'$ rather than left as $16$bis. Page~8 proposes cofibrations injective on objects and arrows with a further call (19), then contrasts general $n$-stacks with $n$-Gr-stacks and their truncated homotopy types. Its wording ``left derived functors'' for the inverse-limit cohomology construction should be treated as an authorial formulation requiring care, not silently changed to the usual right-derived terminology. The final paragraph on p.~9 introduces the dependence on a coherator and the universal-algebra question, despite professing little inclination to pursue it.

\subsection*{17. The Artin--Mazur cohomological criterion (p. 9)}
This short section supplements the unification program by recalling that weak equivalences can be detected through cohomology with twisted coefficients, including noncommutative $H^0$ and $H^1$ as well as the commutative higher groups. Grothendieck explains how this permits étale cohomological information to constrain homotopy types, while wondering whether its success discouraged the development of higher noncommutative cohomology. The reference back to section~15 is the author's own handwritten cross-reference and is retained even though the immediate unification discussion occupies section~16 in the present numbering. The section bridges that motivation and the concrete corrections of section~18. Its handwritten insertion reads ``(namely, unification)'', not the arXiv ``messy unification''; the final sentence deletes ``that'' and says ``their observation'', distinctions that are obscured by the witness. The section number itself is corrected from 18 to 17 on the scan.

\subsection*{18. Corrections to the coherator and higher coherence (pp. 10--11)}
The return to the letter (L) is explicitly corrective. Grothendieck withdraws the existence of arbitrary binary amalgamated sums in $\underline B_o$: the geometric intuition of gluing contractible cells along embedded subcells had concealed the presence of nonmonomorphic maps, including the codiagonal. He also revises the construction of $\underline C_{n+1}$ from $\underline C_n$ by insisting on the total set of new arrows, and rejects the assertion that intersections of maximal subcells must be subcells. The marginal drawing on p.~10 makes the latter failure concrete: three regions $f,f',f''$ meet so that $f\cap f''=\{s,t\}$. The scan therefore supplies a two-point intersection, not the quite different formula in the initial transcription. A separate correction explains why associativity beyond primary composition cannot be represented by a single homotopy arrow: the sources and targets are themselves only homotopic, so squares, cubes and higher coherence data intervene. Page~11 names the dual categories left and right coherators, then connects the proposed non-groupoidal stacks to Benabou's bicategories and their cellular, topological interpretation. This prepares section~19's realization and singular-stack adjunction. The source wording $L\to K$ and $M\to K$ beside $K\amalg_L M$ on p.~10 is retained as an apparent authorial slip; it does not give the usual span for that pushout. Calls (11) and (14) on p.~11 are present without the explanatory footnotes printed in arXiv.

\subsection*{19. A provisional itinerary and the sheaf temptation (pp. 12--13)}
The first two items of the itinerary are the definition of $\underline B_o$, its dual and the coherators, followed by a realization/singular-stack adjunction with topological spaces. Unlike arbitrary presheaves on a left coherator, Gr-stacks must carry the designated amalgamated sums to fiber products. The realization functor nevertheless extends to all presheaves by preserving colimits; its values on basic cells are prescribed, while extensions to the coherator require choices. The anticipated independence of those choices is therefore tied to an as-yet undeveloped notion of equivalence between Gr-stacks. On p.~13 Grothendieck asks whether the exactness conditions could instead be sheaf conditions for a topology, which would make Gr-stacks a topos, and proposes weak equivalences detected by realization, monomorphisms as cofibrations and a lifting-property definition of fibrations. These are programmatic proposals rather than established results. The ``quick glance'' suggesting a topology is retracted immediately in section~20; even the concluding sentence about the inclusion, revised to composite, functor is already struck on p.~13. The source has ``one reward'', not the witness's ``one ward'', and handwritten ``as cofibrations'', not ``for cofibrations''. Its call (13) contains no footnote text on this page.

\subsection*{20. Model categories, sites and a retracted canonical topology (pp. 14--15)}
The section opens by correcting section~19's ``quick glance'': the proposed topology on the coherator would make the component functors fiber functors and force componentwise colimits, which the author recognizes as false. The category of categories supplies a warning and a new candidate construction. Grothendieck considers families of functors whose nerves cover every finite composable string, proposes identifying them with universally effectively epimorphic families, and associates to a category $A$ the topos of sheaves on $(\mathrm{Cat})_{/A}$. The handwritten margin immediately marks the asserted equivalence as false and points to section~24; the digits are readable, while pencil hatching leaves the note's cancellation status uncertain. Page~15 broadens the proposal to a natural site structure on model categories, with the homotopy type of $X$ identified with that of the induced site or topos over $X$. Its further presumption that the topology should be canonical is itself contradicted by a dated margin note absent from the arXiv comparison. These local corrections are essential to reading the section as an evolving investigation. The four arrows in the diagram of effective, universal and ordinary epimorphisms are handwritten, as are the topos tildes in $(M_{/X})^{\sim}=M^{\sim}_{/X}$; the object labels are typed. The inquiry leads into section~21's wider inventory of models rather than establishing a general theorem about canonical sites.

\subsection*{21. Multicomplexes, symmetry and restricted coherence (pp. 15--16)}
The inventory now includes symmetric variants of semi-simplicial and cubical complexes and their multiple-complex analogues. A diagonal complex already describes the intended homotopy type, but Grothendieck explicitly pursues the variety of algebraic presentations, not a unique most convenient model. He compares the passage from complexes to multicomplexes with that from ordinary categories to higher categories or stacks, recalling Quillen's suggestion of defining higher categories through multiple semi-simplicial complexes. The return to the coherator is concrete: the coherence diagrams examined in Benabou's bicategory axioms use restricted amalgamated sums describable through multisimplices, suggesting a smaller starting category $\underline B_o$ than the generous one in the original letter. The handwritten cross-references on p.~16 revise the old page references to sections~8 and~9, and the scan preserves the tense and spelling ``This I ckecked'', where arXiv reads ``This I check''. The suggested nerve-like functor from $n$-stacks to $n$-semi-simplicial complexes is a goal for later investigation. Section~22 next reverses direction, seeking simpler special cases of categories, beginning with preorders and ordered sets. The uncertain final singular/plural in ``connection(s)'' on p.~15 remains an image-edge problem and has no mathematical consequence for this comparison.

\subsection*{22. Ordered sets, realizations and maquettes (pp. 16--19)}
Preorders and ordered sets offer simpler categorical presentations of homotopy types. He relates ordered strata in conically stratified spaces to a realization whose barycentric simplices are finite totally ordered subsets, or flags. The resulting realization is canonically compared with that of the nerve. On p.~17 the corrected chain is $(\mathrm{Ord})\hookrightarrow(\mathrm{Preord})\hookrightarrow(\mathrm{Cat})\hookrightarrow(\mathrm{Ss\ sets})$, not the reversed first pair in arXiv. He cautions that the adjoint passage from a category to its object preorder can lose homotopy information, as the one-object case shows; the handwritten qualification ``but not Fl(A)'' is indispensable to that example. Page~18 gives a second realization through a nonseparated topology on the underlying ordered set and identifies sheaves with covariant functors, hence presheaves on the opposite order. Reversing the order preserves the flags and yields homeomorphic geometric realizations, while the two associated presheaf topoi need not be equivalent. The scan confirms hats in the latter claim and the author's rounded inclusion sign for the initial order comparison. Finally, maquettes are described by vertices and finite simplices and embedded in ordered sets through their face posets. The body claims that this category lacks finite products, while a dated marginal note adds that the notion is not an algebraic structure species; the hatching and call number of that note remain uncertain. Section~23 consequently shifts attention from satisfying Quillen's full package to constructing and localizing a natural functor into $(\mathrm{Hot})$. The edition preserves these claims and reservations as stages of the investigation, without treating the later witness's explanatory footnotes as text written on these pages.

\subsection*{23. The functor to homotopy types and a cellular core (pp. 19--21)}
Grothendieck separates the initial construction of a natural functor $M\to(\mathrm{Hot})$ from the further demand that $M$ carry a Quillen model structure. The functor would define weak equivalences by inverse image of isomorphisms; only then would one choose monomorphisms as cofibrations and test the resulting lifting-property definition of fibrations. The proposed site-theoretic construction uses the homotopy type of the induced site over each object, provided its pro-homotopy type is essentially constant. The canonical-topology suggestion is already qualified by a handwritten retraction pointing to section~24. Page~20 adds a structural requirement, universal and disjoint coproducts, and presents an algebraic species by a category $\underline C$ with finite inverse limits, its dual $\underline B$, and the anticipated equivalence $\mathrm{Ind}(\underline B)\simeq M$. This connects localization to finite presentations and to the question of which exactness properties are already visible in $\underline B$. Page~21 then isolates a smaller category $\underline B_{\mathrm{oo}}$ of basic cells and monomorphic boundary maps: the ordered set of proper subobjects should realize the boundary sphere of the corresponding cell. That construction uses section~22's ordered-set realizations and develops section~18's correction of the initial gluing intuition. The distinction between the broad species-level proposal and the geometric cell data remains heuristic. In the scan, the long p.~20 insertion includes ``can be interpreted'', omitted by arXiv, and the identity-functor arrow is drawn with two shafts; p.~21 preserves the printed slip ``of cubical'' and the typed replacements ``define'', ``family'' and ``sphere''.

\subsection*{24. Truncated nerve topologies and the failure of uniqueness (pp. 21--24)}
This section supplies the correction repeatedly anticipated in the earlier margins. Covering all composable pairs does not cover arbitrary strings: Grothendieck defines a descending sequence $T_N$ by epimorphic families of $N$-truncated nerves, identifying the first two levels with universal epimorphy and universal effective epimorphy. The comparison lemma relates the associated sheaf topoi to presheaves on truncated simplex categories. The finite levels retain only skeletal information, so the full nerve topology $T_\infty$ is needed for the proposed homotopy-type construction. The claim on p.~22 that an $N$-skeleton has the same homotopy and cohomology invariants through dimension $N$ requires caution at the endpoint: the $N$-skeleton of an $(N+1)$-simplex is an $N$-sphere, although the full simplex is contractible. The edition preserves the author's inequality rather than silently strengthening or repairing this intermediate assertion. The same loss of information explains the failure of $K\to NN'(K)$ discussed in section~16$'$: $N'$ factors through the two-dimensional data, while the nerve itself can still induce the desired equivalence after localization. Page~23 then replaces simplicial diagrams by commutative cubical diagrams, producing a strictly coarser topology with a different associated topos but the same intended homotopy interpretation. Multiple complexes suggest infinitely many further suitable topologies; the conclusion is that no unique topology is privileged. This leads directly to section~25's search for an equivalence relation among them. The scan's slice notation on p.~22 places $/A$ outside the parentheses ending after $T_N$, followed by an additional handwritten closing parenthesis; the earlier transcription's placement inside the subscript was corrected, while the extra parenthesis remains as written.

\subsection*{25. Hot-equivalence of topologies and global asphericity (pp. 24--25)}
The lack of a privileged suitable topology leads to a proposed equivalence relation. For comparable topologies $T\geq T'$ on $M$, Grothendieck compares the induced morphism of sheaf topoi and asks that every induced morphism over an object $A$ preserve homotopy type. He calls the resulting property ``globally aspheric'', strengthening the familiar cohomological notion of global acyclicity by allowing noncommutative coefficients. The primitive relation between comparable topologies is not yet an equivalence relation; taking the relation it generates is a first proposal, and he asks whether a coarser useful notion could do better. The intended invariant is a canonically associated functor $M\to(\mathrm{Hot})$, perhaps strong enough to recover the topology class itself. The analogy with Giraud's recovery of a topology from its associated subtopos motivates, rather than proves, a hoped-for correspondence. Page~25 narrows the aspiration to a single natural Hot-equivalence class for suitable species, admitting that even $M=(\mathrm{Cat})$ has no intrinsic characterization yet. Section~26 then reformulates the problem through small subcategories. Two revisions on p.~24 are essential: ``inverse image'' becomes ``direct image'', and the prime is moved so that the inclusion is from $T$-sheaves into $T'$-sheaves. The vertical margin note actually says ``just means'' and ``namely contractible'' with the qualification ``in the Čech sense''; the initial transcription's ``weakly contractible'' was not a literal reading. Its pencil shading leaves cancellation status uncertain. The body ends with the typed slip ``he'' for ``be'', and its final ``this'' is visibly canceled before the continuation on p.~25.

\subsection*{26. Test categories, test functors and the passage back to $M$ (pp. 25--29)}
The proposed intrinsic characterization of suitable topologies is reformulated through small categories $B$ inside a model category $M$. The cellular cores of section~23 and the simplex and cube categories of section~24 suggest that the presheaf category $\widehat B$ should itself describe homotopy types: its natural functor to $(\mathrm{Hot})$, computed from categories of objects over a presheaf, should be a localization. Grothendieck separates this internal requirement on $B$ from the eligibility of a functor $B\to M$; marginal terminology on p.~27 calls them ``test categories'' and ``test functors''. The comparison lemma gives an equivalence between the induced sheaf topos and $\widehat B$ under either of two fullness hypotheses. The decisive qualification comes on p.~28: localization of that sheaf category is not yet localization of $M$ itself. Generation by strict epimorphisms is considered as a way of making $M\to\widehat B$ fully faithful, but the author then proposes dropping both fullness requirements and directly demanding that $M\to\widetilde M\to\widehat B\to(\mathrm{Hot})$ be a localization. For categories $\underline B$ of finite presentations he further hopes for a localization onto homotopy types of finite type. The questions whether equal weak equivalences determine the functor, and whether suitable choices form one natural class, feed sections~27--28. The margin of p.~27 actually says ``same question arises'' for Hot-equivalence of topologies; ``arises'' was absent from the initial transcription. Its wording is now readable, but pencil hatching still leaves the note's cancellation status uncertain. The call after objects ``over a'' on p.~26 is (23), not (3).

\subsection*{27. Geometric realization as a third construction (p. 29)}
Calling the preceding discussion ``thin air conjecturing'', Grothendieck warns that both test categories and test functors may need stronger conditions. He places geometric or ``topologic'' realization beside the two constructions through topologies and through test functors. Semi-simplicial and cubical complexes offer comparatively direct realizations, hence so do the categories embedded in semi-simplicial sets. For Gr-stacks, the prescribed realization of the primitive $\underline B_{\mathrm{oo}}$ structure has to extend to a left coherator, and that extension involves the arbitrary choices already encountered in sections~18--19. Rather than immediately eliminating them, he compares this indeterminacy with the variety of test topologies and functors developed in sections~24--26. The paragraph thus states the next intended problem, relating realization to the two test constructions; section~28 immediately postpones it in favor of rigidity of $(\mathrm{Hot})$. The scan preserves ``topologic'', where the witness expands it to ``topological'', and the word ``is'' in ``is that a suitable'', where the witness has ``it''. Its call (28) supplies no text of the witness's retrospective footnote. The typed insertions ``localization'', ``left'' and ``between'' sharpen the objects of the comparison; the two round letters in $\underline B_{\mathrm{oo}}$ and the primitive graph designation remain literal rather than silently replaced by a joined infinity glyph.

\subsection*{28. Rigidity of $(\mathrm{Hot})$ and the first modelizers (pp. 30--33)}
The intended return to geometric realization is suspended by an ``Assumption'': the category of self-equivalences of $(\mathrm{Hot})$ and their natural isomorphisms should be equivalent to a point. Grothendieck separates its two demands, every self-equivalence is isomorphic to the identity and the identity has no nontrivial automorphism, while wondering whether arbitrary natural endomorphisms could also be admitted. This is asserted as a plausible assumption, not proved; it supplies the desired uniqueness up to unique isomorphism of a localization functor once its weak equivalences are known. A ``modelizing category'' or ``modelizer'' is then a category $M$ with a saturated class $W$ whose localization is equivalent to $(\mathrm{Hot})$. A model-preserving functor carries weak equivalences to weak equivalences and induces an equivalence of localizations, equivalently commuting up to isomorphism with the type functors. This formalizes sections~26--27's questions without yet specifying cofibrations and fibrations. The diagram on p.~31 compares ordered sets, preorders, categories, semi-simplicial sets and cubical sets; the category-of-elements functors $i_A$ lead to right adjoints $j_A$. Page~32 carefully distinguishes those right adjoints from the familiar nerve functors, contrasting infinite categories of simplices with finite ordered face structures when degeneracies are omitted. Mixed multicomplexes then enlarge the diagram and motivate the test-category definition in section~29. The mixed-product display on p.~32 visibly has $\Delta$ in both factors, although the adjacent prose says $p$ times simplicial and $q$ times cubical; this apparent authorial symbol slip is preserved. The same page changes ``modelators'' to ``modelizers'' by hand, a correction the witness does not apply, and contains a small branching diagram omitted there. The p.~30 cross-reference is to section~26; calls (28$'$), (29) and (30) are present without the witness's expanded footnotes.

\subsection*{29. The basic modelizer and a stronger test condition (pp. 33--36)}
The category of small categories now becomes the foundational modelizer: $(\mathrm{Hot})$ is its localization, and weak equivalences are described through the induced morphisms of presheaf topoi and the Artin--Mazur cohomological criterion recalled from section~17. A small category $A$ is provisionally a test category when $i_A\colon\widehat A\to(\mathrm{Cat})$, sending a presheaf to its category of elements, is model preserving; its presheaf category is an ``elementary modelizer''. This replaces a special preference for semi-simplicial sets by a comparison among all such presentations, while acknowledging the computational convenience of the nerve and the simplicial construction of pro-homotopy types. The general adjunction calculation of pp.~34--35 asks when both adjoints induce inverse equivalences after localization. Applied to $i_A\dashv j_A$, it leads to condition (T): every counit $i_Aj_A(B)\to B$ should be a weak equivalence. Grothendieck explicitly leaves open whether this strengthens the preceding definition, and proposes strengthening it if necessary; the call (33) does not contain the witness's retrospective answer. Products of test categories and direct comparison of elementary models become the next questions. The closing heuristic on p.~36 would test (T) by a fibration with contractible fibers, either after taking nerves or directly in $(\mathrm{Cat})$; section~30 immediately attacks the model-category assumption behind the latter route. On p.~35 the source twice writes that $M'$ or $M$ is saturated, although the discussion concerns the arrow classes $W'$ and $W$; these apparent symbol slips remain unchanged. The typed ``elementary normalizer'' on that page likewise remains source text. The scan's baseline slash in $M/_{b}$ on p.~34 is preserved, and its reference reads ``section 28'', not the witness's page reference. The broad emphasis of the witness also differs from the selective underlining and marginal double bar of p.~33.

\subsection*{30. From the model-category obstruction to asphericity (pp. 36--38)}
The attempt to treat every monomorphism of categories as a cofibration now runs into the two elementary horns drawn on p.~37. Lifting against both would force the arrows of a fibrant category to have left and right inverses, leaving groupoids, whose homotopy types cannot exhaust those of arbitrary categories. The objection concerns that proposed choice of cofibrations; the source does not prove the nonexistence of every possible model structure on categories. Grothendieck nevertheless keeps the expectation that elementary modelizers should admit closed model structures, and turns instead to the asphericity criterion introduced in section~25. A morphism of topoi is required to remain a weak equivalence after every localization over the target; when the topoi come from categories, generators reduce the condition to the asphericity of each slice $C'_{/a}$. Applying this to the counit $i_Aj_A(B)\to B$ from section~29 gives three equivalent conditions: the strengthened test condition, asphericity of every counit, and asphericity of $i_Aj_A(B)$ whenever $B$ has a final element. That last formulation opens the explicit interval argument of section~31. The scan's four bracketed ``[small]'' restrictions on p.~38 are genuine hand additions, including two absent from the witness. Its caret under the $A$ in $\operatorname{Hom}(A_{/a},B)$ is retained in its actual position without asserting a mathematical interpretation. The apparent inline hat on $f$ reported by the initial comparison on p.~37 belongs instead to the display beneath it; the inline symbol has been corrected to plain $f$. The word ``closed'' on p.~37 and the replacement ``aspheric'' on p.~38 are interlinear typing, although their placing strokes are handwritten.

\subsection*{31. An interval proof of the test-category criterion (pp. 38--41)}
The search for a direct proof that the simplex category is a test category yields three sufficient conditions. T1 asks that $A$ be aspheric; T2 asks that $A_{/a\mathrm{x}b}$ be aspheric for any two objects, with their product interpreted in the presheaf topos when necessary; T3 supplies an aspheric interval object $I$ with disjoint final subobjects $e_o,e_1$. For $F=j_A(B)$ with $B$ admitting a final element, the aim is a homotopy from the identity to the constant section. Diagram (D) on p.~39 expresses the cylinder $F\mathrm{x}I$ and its two endpoint inclusions. Given $u:a\to I$, the inverse image of $e_o$ determines a sieve $C_u$ in $A_{/a}$; a functor $f$ is extended by its own values on $C_u$ and by the final object of $B$ on the complementary objects. The sieve property makes the remaining arrows uniquely definable. This constructs the homotopy without T1, T2 or the asphericity of $I$. Those hypotheses enter the second step: the slice criterion of section~30 reduces the asphericity of $F\mathrm{x}I\to F$ first to $a\mathrm{x}I$ and then to $a\mathrm{x}b$, precisely T2. T1 then converts the weak equivalence to the final presheaf into absolute asphericity. Page~40 separately justifies the retract step through the Artin--Mazur criterion rather than claiming it is formal for an arbitrary arrow class. The scan gives genuinely hooked tails to both arrows into the cylinder in (D), contrary to the initial transcription; all diagram arrows are handwritten and its object strings are typed. The p.~39 reference is ``cf théorème p 46'', while the marginal comparison involving T1 and T2 retains an uncertain short connector. The live punctuation after ``checked'' on p.~41 is a colon, and the source's ``asperic'' remains uncorrected.

\subsection*{32. A revised programme after the interval criterion (p. 42)}
Dated 15.3, the new planning page treats T1--T3 as convincing candidates rather than a finished characterization. Its first three tasks are to examine necessity, construct the familiar simplicial and cubical examples, and prove stability under products so that different elementary models can be compared. The fourth returns to section~30's unresolved closed-model question; the fifth restores the larger programme of modelizers attached to algebraic structure species and the possible uniqueness of their weak equivalences. The author explicitly notices that stacks have receded into the background. A sixth task broadens the inquiry to a satisfactory replacement for Verdier's provisional triangulated-category notion, including the exactness, generator and cogenerator questions for $(\mathrm{Hot})$. Section~33 begins carrying out the first task and introduces the distinction between test and strict test categories. The scan matters to the direction of the programme: item 5 says categories $M$, not the witness's $\widehat A$. Its call (9) is present without the witness's explanatory footnote. ``Gave to this notion'' and ``general'' are typed interlinear revisions, while ``ways of'' and the displaced ``of'' are handwritten. The cancelled ``topos'' in item 4 is legible beneath both machine and pen deletion; the cropped completion of ``which'' and the old word after ``constructing'' remain doubtful.

\subsection*{33. Strict test categories, the Lawvere interval and a qualified theorem (pp. 43--47)}
The provisional terminology first separates test categories from ``strict'' test categories satisfying T1--T3. T1 is necessary by the final-category case of the criterion on p.~38. T3 admits a canonical candidate: applying $j_A$ to the two-object ordered interval produces a presheaf whose sections are sieves, hence the ``Lawvere element'' representing subobjects. This ties the explicit interval construction of section~31 to an intrinsic object of the presheaf topos. A proposed concluding statement about an intrinsic T3$'$ is completely cancelled on p.~43. The remaining issue is T2. Pages~44--45 connect it with preservation of binary products by the localization functor: the comparison $A_{/F\mathrm{x}G}\to A_{/F}\mathrm{x}A_{/G}$ is a weak equivalence for all $F,G$ exactly when the representable-product condition holds, and the slice criterion in fact supplies asphericity. The blanket claim about finite direct and inverse limits in $(\mathrm{Hot})$ on p.~44 is stronger than the finite sums and products used here; its marginal (!35) must not be expanded into an unqualified theorem in the guide. The result assembled on p.~46 compares product-preserving localization, the localized adjunction $i_A\dashv j_A$, and T1--T3. Its later margin dated 9.4 is essential: condition (i) is said to be ``not strong enough'', and the left continuation proposes adding that $j_A=i_A^*$ is model preserving, with a reference to remark~3 on p.~174. The left note is readable but crossed by broad pencil hatching, so its cancellation status remains uncertain. The source says ``model preserving'', not the witness's ``equally modelizing''. The author also explains how to avoid assuming that $(\mathrm{Cat})\to(\mathrm{Hot})$ preserves binary products: formulate the condition directly for $i_A$ up to weak equivalence. Page~47 provisionally reserves ``test category'' for the strict case, with call (38), and characterizes elementary modelizers through presheaf-topos structure, localization and finite products; a wider modelizing-topos notion is deferred. The scan retains ``asphéric'' on p.~43, three rounded intersection caps rather than the transcription's wedges, typed interlinear additions to the theorem, and the actual oval marginal annotations. The two opposite bent arrows on p.~47 remain graphic annotations, without an invented reordering instruction.

\subsection*{34. Product-stable examples and the homotopy circularity (pp. 47--50)}
The first examples strengthen the sufficient hypotheses again. A final object gives T1$'$, and an actual object $I$ of $A$ with two disjoint sections gives T3$'$; if $A$ is also stable under binary products, T2 follows because each product is representable and its slice has a final object. This constructs test categories inside any suitable ambient category, even one that is not small, by selecting a small product-stable subcategory. The economical construction uses powers of an interval and only projections and constant endpoint maps as components. Its cubical category includes permutations and diagonals, so p.~49 distinguishes it from the more familiar semicubical category. Nonempty finite sets give a corresponding full simplicial construction, with affine simplices as its geometric interpretation. The notation on p.~48 literally says to include $I$ and $\Delta_o,\Delta_1$; these triangular symbols are retained even though the preceding discussion specifies endpoint maps $\delta_o,\delta_1$. Returning to the ordinary simplex and cube categories removes the convenient product-closure hypothesis. Grothendieck therefore tries to prove T2 by making every representable deform to the final presheaf. The decisive revision on p.~50 is his own detection of a circular argument: section~31's homotopy proof used $F\mathrm{x}I\to F$ as a weak equivalence, and that was established from T2. A separate condition, called (2), that every $a\mathrm{x}I$ be aspheric must therefore accompany the deformation-retract condition TH1. Section~35 postpones this unfinished verification for the intrinsic theory of modelizing topoi. The formula $u_a=p_a c_a$ on p.~50 is retained in that visible order; with $c_a:e\to a$ and $p_a:a\to e$, the surrounding right-to-left composition convention would instead require $c_a p_a$ for an endomorphism of $a$. The scan also confirms the literal spellings ``axists'', ``oddurs'' and ``consequance''. The extra presheaf caret on the full simplicial category on p.~49 is present despite the witness omission, and the stepped mark on p.~48 is retained graphically in the diplomatic reading and applied as a paragraph break only in the master reading.

\subsection*{35. Local and total asphericity of topoi (pp. 50--55)}
The detour from section~34's unfinished simplex proof translates its conditions into the language of topoi. An aspheric topos is required to be $0$-connected and to have trivial cohomology with constant coefficients in every defined positive degree, including the noncommutative first degree. For an object or arrow, the induced topos supplies the corresponding notion. A locally aspheric topos has a generating family of aspheric objects; these generators suffice for testing whether an arrow remains a weak equivalence after every localization (p.~51). Proposition~1 identifies product stability of aspheric objects with their being aspheric over the final object. Total asphericity combines that property with local asphericity, and the handwritten ``non empty'' inserted on p.~52 is essential: without it, the preceding assertion that T2 already implies T1 would admit the empty case. The ensuing \v{C}ech argument gives absolute asphericity. For topological spaces the much stronger total condition becomes irreducibility, so a Hausdorff example must be a point. This sharpens section~33's distinction between a local supply of cells and the global product condition. Proposition~2 then separates the interval requirement from total asphericity. Its original ``totally aspheric'' hypothesis is visibly cancelled, and p.~53 explicitly says that no local-asphericity assumption is being made. The remaining equivalence is between the Lawvere element being aspheric over the final object and the existence of an object with that property and two disjoint sections. The homotopy argument of section~31 supplies the converse, and localization upgrades the resulting weak equivalence to asphericity. A connected topological space provides a useful obstruction: the proof with direct summands of its Lawvere element shows that this element is disconnected precisely in the irreducible case, where it has two connected components. Consequently no topological space can simultaneously meet total asphericity and the interval requirement (pp.~53--54). The proposed term ``modelizing topos'' packages both, while a second, still conjectural characterization asks for a saturated localization equivalent to $(\mathrm{Hot})$, compatible with finite sums and products and with the Verdier--Artin--Mazur construction. The closing inquiry about morphisms $X\to\operatorname{Top}(B)$ seeks a categorical description closer to the given indexing category than a large inverse system of semi-simplicial objects (pp.~54--55). The source retains ``0-connexity'', ``independant'', ``Corollay'' and ``monomorphisme''. Its upper margin on p.~54 reads ``cf (34)'' beside a body call (39); the lower strictness-related note remains partly doubtful. These are retained source features, not expanded into the arXiv witness's retrospective commentary.

\subsection*{\S67. Aspheric maps replace contractibility as the organizing notion (pp.\ 188--191).}
Grothendieck reorients the modelizer program: the crucial objects are aspheric objects and weak equivalences, and its natural comparison maps are aspheric functors. He explicitly judges the earlier restriction to contractibility and special test functors to have been an unnecessarily long detour (p.\ 188), while retaining its possible relevance for ``canonical modelizers.'' The discussion now allows arbitrary small categories, rather than only test categories. For a basic localizer $\underline W$ he introduces $(\mathrm{Hot}^{W}_{A})$, also written $(\mathrm{Hot}_{A})$, as the localization of the presheaf category by $\underline W_A$; the construction is then organized by inverse image $i^*$ and the comparison morphism $\lambda_i(F):A_{/i^*F}\longrightarrow B_{/F}$, interpreted through the corresponding induced topoi (pp.\ 189--190). This prepares the promised description of asphericity structures and testing functors at the end of the section (p.\ 191); the next section interrupts that development. The dated marginal correction of 3 May on p.\ 189 expressly retracts the claim that Quillen's weak equivalences are precisely $\underline W_A$, saying that Quillen's are stronger. It is therefore essential to read the original assertion together with the margin. The witness changes handwritten ``or'' into equality in (1), drops the subscript $A$ on $\underline W_A$, omits the handwritten domain and codomain after the action on $(a,p)$ on p.\ 190, and prints the struck word ``real'' on p.\ 188. These differences concern notation, functorial context, or the author's revision; none licenses an unrecorded mathematical repair of the typescript.

\subsection*{\S68. Mathematical work, the Teichmüller tower, and the necessity of groupoids (pp.\ 191--195).}
Dated 27 May, the section explains the six-week interruption in the asphericity discussion through a retrospective of the author's changing mathematical interests. Its mathematical center is the ``Lego-Teichmüller construction game'': cutting and gluing surfaces, the tower of groupoids $T_{g,\nu}$, and the arithmetic action on their profinite completions (pp.\ 191--192, 194--195). The account of \emph{La Longue Marche}, tame topology and ``unfolding'' of stratified structures supplies the context for the anabelian project; the intervening reflection on creative interest, ambition and alternating periods of work belongs to the argument's historical presentation and must not be omitted as an aside (p.\ 193). The final page gives the technical lesson: reducing several natural base points to one can destroy visible symmetries and obscure a presentation by generators and relations, whereas groupoids preserve the geometry. Van Kampen situations supply the comparison, and the modular ``multiplicities'' are identified as topoi whose points have finite automorphism groups (p.\ 195). This is a return to the text's groupoid motivation, rather than a proof of the asphericity claims left pending in \S67; \S69 resumes a retrospective of the technical scratchwork. The section explicitly revises an earlier preference for groups or ``exterior groups'' in favor of groupoids (p.\ 195). The witness has several ordinary wording errors here---including ``case'' for ``care'' on p.\ 195---while the underlining and handwritten exclamation in ``topoi!'' preserve the emphasis of the typescript.

\subsection*{\S69. Derivators and integration of homotopy types (pp.\ 196--202).}
The technical retrospective explains why an ordinary derived or homotopy category retains too little structure: traces, determinants and the compatibility of iterated constructions motivate the system of localized diagram categories $D_A(I)$, contravariant in the indexing category $I$, with inverse-image functors expected to admit both adjoints (pp.\ 196--198). This develops the modelizer program recalled in \S67 and connects it with the geometric motivations of \S68. For categories, a covariant diagram can be integrated by forming its associated cofibered category; the comparison with localized slice categories yields the model for ``integration of homotopy types'' (pp.\ 198--199). The name ``derivator'' is introduced on p.\ 200, together with the internal stability expressed by $D(I\times J)$ and the aspiration to a corresponding homotopy theory of derivators. The geometric example is the system of tubes around strata of Mumford--Deligne compactifications: their $K(\pi,1)$ homotopy types, with Teichmüller-type groups, suggest an algebraic description from a diagram of fundamental groupoids (p.\ 201). A diagram in $(\mathrm{Hot})$ alone is explicitly said to be insufficient; a diagram in a model category is required. Grothendieck retracts his recent expectation that finite limits exist in $(\mathrm{Hot})$ (p.\ 201), and \S70 immediately examines the extra factorization condition proposed for the two adjoints (pp.\ 200, 202). Significant witness differences include $T$ replaced by $\mathscr D$ in the statement about inverse limits (p.\ 198), $q=f\circ p$ replaced by an equation beginning with $g$, and ``basic modelizer'' replaced by ``basic localizer'' (p.\ 199). These readings are clear on the images and are preserved rather than repaired from the witness.

\subsection*{\S70. Base change, duality, and the limitations of $(\mathrm{Cat})$ (pp.\ 202--206).}
Starting from the desired two-stage factorization $X\to Z\to Y$, Grothendieck asks for preservation of weak equivalences by base change along the second map and co-base change along the first, with either factor optionally a weak equivalence (p.\ 202). The goal supplies the extra condition proposed for derivator operations in \S69. The investigation of the base-change side leads to cohomological smoothness and properness of functors, characterized through asphericity; fibered categories are smooth, cofibered categories proper, and passage to opposite categories interchanges the two conditions (p.\ 203). The pairing of presheaves and copresheaves suggests a broader duality of topoi, but the claim that cosheaves on an arbitrary topos form a topos remains a question (pp.\ 203--204). The co-base-change side is less tractable: open or closed immersions in $(\mathrm{Cat})$ fail the hoped-for general assertion, whereas pushouts along monomorphisms in $A^{\wedge}$ have the required homotopy meaning through the integral of the associated span (pp.\ 204--205). The stated smooth/proper exceptions and the possible special case of an immersion that is also a weak equivalence must be distinguished from that rejected general assertion. This forces a return to elementary modelizers and Quillen factorization questions; the author finally accepts the Kan condition after describing his resistance to it, while retaining dissatisfaction with the reduction to semisimplicial methods (p.\ 206). These questions become part B of the program in \S71. The witness omits the monomorphism hook in the p.\ 204 square and replaces the actual objects $Y,X,Y'$ of the p.\ 205 integral diagram by categories of elements. The scan-confirmed diagrams preserve the distinction between these two levels of the construction.

\subsection*{\S71. A provisional program: modelizers, derivators, and topoi (pp.\ 206--212).}
The program dated 30 May organizes the preceding scratchwork into four tasks. Part A returns to the asphericity structures and canonical modelizers suspended in \S67; part B develops the base-change, co-base-change and factorization questions of \S70, with homotopy integrals, amalgamated sums and Mayer--Vietoris sequences as the connecting constructions (pp.\ 206--207). A handwritten transposition on p.\ 207 makes strong saturation one of the properties expected to follow from the first three basic-localizer axioms. Part C develops the derivator proposal of \S69 while asking what additional structure is needed for external Hom's and basic invariants. It then distinguishes group objects modeling pointed connected homotopy types from the truncated output of group objects in $(\mathrm{Cat})$, and asks which modelizers satisfy a Dold--Puppe theorem (pp.\ 207--209). The proposed trilogy consists of semisimplicial, cubical and hemispherical complexes, the last also called globular complexes by Brown; their boundary and degeneracy operations are related to regular cellular subdivisions of spheres. Multicomplexes give multiple chain complexes, so the simple Dold--Puppe statement does not extend unchanged. The classification of the three series and the need for enough degeneracy maps are presented as questions to investigate, not proved theorems (pp.\ 209--210). Part D, ``Back to topoi,'' seeks relative homotopy categories $D(T)$ and their possible opposite counterparts, extending the $f^*,f_!,f_*$ formalism from indexed categories to topoi. The discussion keeps existence of derived adjoints tentative and moves from objects to proobjects when a topos fails to determine an essentially constant homotopy type (pp.\ 210--211). Local asphericity then reconnects the program with \S35; the revised distinction between modelizing and strictly modelizing topoi is recalled from p.\ 68, and the section ends with the Lawvere-element condition for a finitely triangulated contractible space (p.\ 212). The next section begins the promised return to the asphericity game. The witness changes $f_*$ to $f^*$ and $Lf_!$ to a subscript-$1$ expression on p.\ 210, and applies a handwritten transposition there without marking it. The duplicated target-only display on p.\ 211 was an error of the draft transcription, absent from the scan; its removal restores the actual three-display sequence.

\subsection*{\S72. Criteria for aspheric functors and the first stability results (pp.\ 213--219).}
Returning to \S67 after the program of \S71, Grothendieck isolates the localizer assumptions needed for asphericity: mild saturation, the condition on categories with final object, and the locality condition Loc~3 (pp.\ 213--214). Proposition~1 identifies the comparison maps on categories of elements, their restrictions over objects, and the corresponding slice criteria; it then distinguishes these equivalent asphericity conditions from preservation of weak equivalences and equivalence after localization (pp.\ 214--216). The extra hypotheses in the implication diagram matter: the discrete-target example separates aspheric fibers from merely nonempty fibers, and the product-projection example shows that inducing the localized equivalence need not determine an aspheric map without the stated assumptions (p.\ 216). Under asphericity, the localized triangle commutes up to the canonical comparison isomorphism; this reconnects the discussion with the earlier ``inspiring assumption'' of \S28 without using that assumption as a proof (pp.\ 216--217). Proposition~2 transfers total asphericity along an aspheric functor and relates the inverse image of a separating homotopy interval to asphericity of the functor. Its handwritten addition makes the source category aspheric, and the long interlinear insertion specifies the homotopy condition (TH1) of p.\ 50; the corollary deduces strict test categories from a totally aspheric source and a local test target (pp.\ 217--218). Proposition~3 supplies composition and right cancellation, base change along a fibering map, and products (pp.\ 218--219). The base-change proof uses the fibration condition L5, whose derivation from the stated localizer assumptions is expressly postponed; Grothendieck explains that only part~(a) is needed for the immediate asphericity-structure formalism, preventing the announced application from relying on that deferred argument (p.\ 219). Remark~4 denies the opposite cancellation assertion and leads directly to the examples in \S73. The visual collation restores the proof-diagram arrow from (i'') to (iii), rather than to (i'), on p.\ 215; it also identifies the apparent deletion of the t in Hot on p.\ 214 as pen retracing. On p.\ 218 the comparison sign is a handwritten tilde over a bar, with no arrowhead, contrary to the draft transcription and its diff description. These are distinctions between what the image shows and what a mathematically plausible reconstruction might supply.

\subsection*{\S73. Adjoints, coasphericity, and cofinality (pp.\ 220--221).}
The continuation dated 11 June begins by suppressing the qualifying $\underline W$ when one basic localizer is understood. It sharpens Remark~4 of \S72: a functor from the one-point category into an aspheric category is aspheric exactly when its chosen object is initial, illustrating how much stronger a map condition is than asphericity of its source and target (p.\ 220). Proposition~4 enlarges the preceding isomorphism example: every equivalence is aspheric, and more generally a left adjoint is aspheric because its slice categories identify with slices having final objects. Passage to opposite categories then motivates the name ``coaspheric'' for the dual condition, satisfied by the right adjoint. The corresponding undercategories and their initial objects connect coasphericity with a promised all-dimensional cofinality statement, whose dimension-zero consequence is invariance of direct limits (pp.\ 220--221). The analogy from \S70 returns: asphericity resembles properness with aspheric fibers, and coasphericity resembles smoothness with aspheric fibers. The comparison is explicitly weakened because arbitrary base change does not preserve either condition; the applicable base changes are fibering/smooth or cofibering/proper, respectively. Even an equivalence of categories can have an empty fiber, so its simultaneous asphericity and coasphericity do not supply arbitrary base-change stability. Universal weak equivalences, called ``trivial Serre fibrations'' here, are proposed as a stronger common condition and assigned to a later systematic investigation (p.\ 221). The section then turns to the overdue product discussion of \S74. On p.\ 220 the crossed handwritten heading above Indeed is Remark~5, not the draft's tentative Proof; it is written again at the subsequent paragraph. Whether the placement strokes also delete its typed initial The remains unresolved and is marked as such. The p.\ 221 witness changes singular condition to plural conditions and supplies a title for \S74 absent from the typescript.

\subsection*{\S74. Products, diagonal functors, and canonical comparison of models (pp.\ 221--230).}
The overdue product discussion supplies concrete constructions for the abstract stability results of \S72. The external product of presheaves has a category of elements identified with the product of the two categories of elements; it preserves aspheric objects and weak equivalences in the stated variables, and yields the product assertions for locally aspheric and totally aspheric categories (pp.\ 221--222). The corollaries distinguish local test, test, and strict test categories: the hypothesis on the other factor changes from arbitrary to aspheric to totally aspheric. Mixed multicomplexes illustrate the construction (p.\ 223), reconnecting it with the trilogy in \S71. The comparison functor associated with an aspheric map between modelizing categories leads to a functor from $\mathrm{Asph}(A,B)$ into a category of comparison functors and factors through its fundamental groupoid. Nonemptiness and connectedness suffice for the existence of an isomorphism class of comparison, while a canonical transitive system requires the stronger simple-connectedness conclusion; the text carefully separates these claims (pp.\ 224--225). The diagrams on p.\ 226 distinguish categories, presheaf modelizers, and localized homotopy categories, including the middle presheaf triangle omitted from the witness.
The diagonal criterion then identifies total asphericity of $A$ with asphericity of $A\to A\times A$, and permits simultaneous comparison through products of strict test categories (p.\ 227). Proposition~4 treats pointwise products of functors into a category $B$ with binary products and aspheric presheaves $\underline{\mathrm{Hom}}(b,y)$: the total-asphericity assumption on $A$, moved by a handwritten transposition into the hypotheses, is essential to the stated construction (p.\ 228). The external-product corollary is sharpened by revision to an equivalence: when $i$ is aspheric, $i\boxtimes i'$ is aspheric if and only if $A'$ is aspheric. Its proof uses a lemma on a presheaf over $A\times A'$ whose restrictions to $A$ are aspheric. The resulting category-of-elements projection is a weak equivalence by the fibration condition L5, explicitly recalling the claim and deferred justification on p.\ 219; this dependency should remain visible (p.\ 229). Contractibility of the objects of $B$ is discussed as a stronger practical hypothesis. Finally, closure of $\mathrm{Asph}(A,B)$ under binary products makes its nonempty case connected with the needed homotopy strength, giving the canonical transitive system of isomorphisms between the localized functors and the canonical functor announced on p.\ 225 (p.\ 230). A handwritten margin points forward to the generalization in \S82, p.\ 272. The image-based edition retains the unusual $X\longrightarrow B_{/X}$ display on p.\ 224, the missing slash in the final category-of-elements notation on p.\ 226, and the visibly typed product x and lowercase subscript o on p.\ 228. The hemisphere-like symbol on pp.\ 223 and 229 remains a doubtful approximate glyph; neither mathematical expectation nor the witness settles its exact drawn form.

\subsection*{\S75. Asphericity structures, testing functors, and weak equivalences (pp.\ 230--237).}
The generalities of \S\S72--74 now become a structure on a possibly large $\underline U$-category $M$. After making the universe convention explicit, Grothendieck defines a class $M_{as}$ of aspheric objects through a small category $A$ and a functor $i:A\to M$ that lands in $M_{as}$ and detects it by the presheaf $i^*(x)=\mathrm{Hom}(i(-),x)$ (pp.\ 230--231). Such an $i$ is a “testing functor.” Every small full subcategory supplies an asphericity structure, and Proposition~1 plus its corollary shows that any structure can be generated this way; enlarging the generating full subcategory within $M_{as}$ preserves the structure (p.\ 232). The following terminology discussion distinguishes testing functors from the earlier test functors of pp.\ 175--176. The essential condition depends on the possibly large category $M_{as}$ itself, rather than on its embedding in $M$, so the author tentatively adopts $M_{as}$-aspheric, or simply aspheric, as a synonym. His use of both “aspherical” and “aspheric” in this discussion is preserved (p.\ 233).
Proposition~2 defines $W=(i^*)^{-1}(\underline W_A)$ and proves independence of the testing functor by passing to a common small full subcategory. Mild saturation is immediate; strong saturation remains conditional on the corresponding result for the basic localizer (p.\ 234). When $M$ has a final object, calling the structure aspheric means that this final object belongs to $M_{as}$; the resulting criterion recovers aspheric objects from the maps to the final object. The small-generator descriptions then characterize both the admissible weak-equivalence sets and the asphericity classes. Without the aspheric final-object hypothesis, a weak-equivalence set need not determine a unique asphericity structure. The analogy with contractibility and homotopism structures in \S\S51--52 motivates a promised construction from contractibility to asphericity and a uniqueness assertion for the intended canonical examples (p.\ 235). Proposition~4 recasts the criteria of p.\ 214 as six conditions: three about aspheric objects and three about weak equivalences. The object conditions are equivalent and imply the others; the revised converse requires both $A$ and the asphericity structure to be aspheric. All six conditions become equivalent when the specified generating category $B$ contains the final object. If no $B$ is fixed, enlarging one to contain $i(A)$ and the final object yields equivalence of the first five conditions (pp.\ 236--237). These added hypotheses, some typed interlinearly and some handwritten, are central to the proposition and were partly concealed by malformed insertion macros in the draft. Section~76 immediately tests the formalism on examples.
Two visible mathematical slips are retained for the reader to assess: the universe paragraph says “two objects of $U$” where the witness substitutes $M$ (p.\ 230), and Proposition~2 says “Then $M$ is” a mildly saturated subset of $\mathrm{Fl}(M)$ after defining $W$, where the witness prints $W$ (p.\ 234). Proposition~3 and its corollary also call the relevant map “aspheric,” while the explicit revisions of Proposition~4 replace that word by “a weak equivalence” (pp.\ 234, 236); the edition does not silently harmonize the earlier wording. Several correction layers remain uncertain, including the usual-localizer subscript at the preceding section boundary and overtyped letters in these definitions and proofs.

\subsection*{\S76. Basic examples and total asphericity (pp.\ 237--239).}
Three examples test the definitions of \S75. For $M=(\mathrm{Cat})$, the $\underline W$-aspheric categories form the basic asphericity structure, generated through $a\mapsto A_{/a}$ for a weak test category $A$. The shift from the basic modelizer to this structure puts aspheric objects before weak-equivalence arrows. The distinction matters for functors into $(\mathrm{Cat})$: asphericity supplies the induced localized functor, whereas the weak-test requirement asks for a homotopy-model comparison. Theorem~1 of p.\ 176 suggests equivalence of these conditions when the values are contractible, and the author identifies a comprehensive answer to the “silly question” of \S46, p.\ 95, as the motivation for the present formalism (pp.\ 237--238). The second example takes topological spaces with the weakly contractible spaces as aspheric objects; the geometric realization of simplices supplies the testing functor, and the text appeals to Gabriel--Zisman for the stronger modelizer statement. Grothendieck adds the previously omitted assertion that the first two structures are aspheric and later notes that the resulting weak equivalences are the usual ones (p.\ 238). The third example is $A^\wedge$ with its usual aspheric presheaves, generated by the canonical fully faithful inclusion of any small $A$. Its weak equivalences are $\underline W_A$, while its asphericity is equivalent to asphericity of $A$; the modelizing property is not asserted for arbitrary $A$. A functor $A'\to A$ is aspheric precisely when its composite into this presheaf structure is $M_{as}$-aspheric (pp.\ 238--239). This leads to “totally aspheric” structures: $M$ admits finite products and finite products of aspheric objects, including the empty product, are aspheric. Proposition~5 identifies this condition with the existence of a small generating subcategory stable under finite products, and with a testing functor from a category whose presheaf category is totally aspheric. It reconnects the product arguments of \S74 with the general formalism and prepares the localized canonical functor of \S77 (p.\ 239). The visual check restores a missing typed a before the repaired l in “realization,” retains the literal thatnthis and appheric typing, and distinguishes typewritten interlinear additions from pen text. The upward margin arrow at the end of the examples had been misrepresented by three brackets in the draft; its direction is restored, without attributing a mathematical meaning to the layout mark.

\subsection*{110. Schematization through Postnikov fibrations, and its obstruction (pp. 443--445; 134-7/1--3).}
Grothendieck resumes the question of homotopy types over a ground ring by comparing an abelian group $M$ with the functor $W(M)(k)=M\otimes_{\underline Z}k$: the discrete object acquires a structure analogous to a vector bundle, suggesting an analogous construction for semisimplicial models and successive Postnikov fibers. The proposed ``schematization above level n'' would leave a lower truncation discrete and require the comparison from schematic cohomology with quasi-coherent coefficients to ordinary cohomology to be bijective. The reference to section 94 situates the hope of replacing $\underline Z$ by $\underline Q$ when nilpotence replaces an abelian fundamental group. On p.~444 he immediately questions this construction: torsion and the universal coefficients formula prevent the homotopy groups after base change from being obtained by tensoring the individual homotopy groups. The alternative is to reconstruct a homotopy type from a suitably enriched ``abelianization'' or ``linearization''; section 111 then seeks a different construction without explicit Postnikov decomposition. The handwritten replacement of ``(= homotopy)'' by ``(or rather, homotopy)'' on p.~444 matters here: the arXiv witness retains the cancelled equality wording. Its mapsto in the displayed base-change formula also differs from the scan's ordinary arrow; the scan alternates star subscripts with the typed notation X. Page 445 closes with a methodological reflection on the tension between detailed preparatory work and the desire to move quickly. The rust-obscured words on p.~443 remain marked rather than restored from the mathematical argument.

\subsection*{111. Unipotent bundles, divided powers, and a revised notion of equivalence (pp. 446--459; 134-7/4--17).}
The starting problem is to construct categories $Hot_n(k)$ of schematized $n$-connected homotopy types, with extension and restriction of the ground ring and a comparison with the ordinary homotopy category over $\underline Z$ (pp.~446--449). Grothendieck builds $U(k)$ inductively from affine bundles $W(M)$ and torsors, first in the fpqc setting. He asks whether the successive fibration data belong to the object or merely witness its existence, then proposes semisimplicial objects subject to connectivity and filtration conditions. The proposed weak equivalences are initially detected by taking sections, and the hoped-for $k$-module structures on homotopy groups lead to a search for a linearization into $D.(Ab_k)$. On p.~448 the model-category symbol has a typed subscript $m$, although the surrounding development uses $n$; this apparent index slip is retained in the text. Pages 450--455 test divided powers as the required linearization. The candidate universal map $x\mapsto\sum x^{(i)}$ immediately forces a completion, and the completed divided-power coalgebra is considered through its group-like elements. After the Eilenberg--Maclane test, the discussion returns to pointed objects and uncompleted divided powers. The comparison (w) on p.~454 places ordinary free abelian linearization, $L$, and the pointed version $L_{pt}$ in one diagram. The marginal correction on p.~452 is consequential: ``One has to take k=Z, and components $X_n$ projective.'' Both its initial words and the subscript are legible on the scan. The subsequent distinction between quasi-isomorphisms and chain homotopisms supplies the projectivity restriction. The final reference in the sentence on p.~455 is to (x''), after a preceding reference to (x'''); the arXiv witness repeats (x''') and loses this distinction. The argument is substantially revised on pp.~456--459. The author retracts the unrestricted derived-functor claim of p.~454, requires projective or flat components, and remarks that the Postnikov approach must again be set aside. The plain typed $I_k(M_*')$ on p.~456 is retained as printed, although the neighboring formulas suggest a notation slip. Page 457 then exposes the problem with ordinary homology as a $k$-module: the exterior-power calculation for $K(M,1)$, applied to $k=\underline Z[T]/(T^2+1)$, yields an obstruction. Page 458 accordingly defines total linearized homology $LH.(X_*)=L_{pt}(X_*)$ and considers weak equivalences detected by its homology invariants. The handwritten revision says it is becoming \emph{dubious} whether equivalences should be defined through $X_*(k)$; the arXiv witness reads \emph{obvious}, reversing the author's change of judgment. At the end of the section, Question 2 asks whether the two definitions coincide for projective pointed unipotent bundles satisfying $X_0=X_1=e$ (p.~459). The weaker condition $X_0=e$ raises an additional issue about abelian fundamental groups. These pages formulate questions and revised expectations, rather than establish an unrestricted equivalence theorem; section 112 resumes from the projective $\underline Z$ case.

\subsection*{112. Kan conditions and the scope of the linearization program (pp. 460--465; 134-7/18--23).}
The entry dated 26.8 takes the projective $\underline Z$ case of the comparison maps (w) as its working foothold. Grothendieck sketches a passage from the Eilenberg--Maclane case to finite chain complexes, then through finite and ``way-out'' infinite successive fibrations, using multiplicativity and hoped-for spectral-sequence or derived-category transitivity statements (pp.~460--461). The proposed Kan--Postnikov complexes use shifted projective resolutions of homotopy groups. To transfer the construction from discrete to schematic objects, he returns explicitly to the cohomology obstruction of p.~444 and hopes to dualize the homology comparison. This is an outline requiring substantial verification, as the author himself stresses. Constant complexes then show why $X_0=e$ is necessary, and the monoid discussion turns the issue into comparison of polynomial with arbitrary cochains. The multiplicative-monoid example on p.~462 is offered as a counterexample to the assumptions previously used; the circled marginal cross is part of the page's evidence and is absent from the arXiv witness. The response is to ``drop Postnikow and keep Kan'' (p.~462). Since $U(k)$ is not closed under arbitrary fiber products, a schematic horn condition must ensure both representability of the relevant finite limits and a suitably strong surjectivity property (p.~463). The tentative name ``submersions'' covers alternatives ranging from trivial bundles to formally smooth maps admitting sections. The first nontrivial condition asks for the three maps $X_2\to X_1\times X_1$ to be submersions; for a monoid nerve this is meant to force a group. Page 464 lists the desired consequences: the linearization theorem over $\underline Z$, a homology criterion for homotopisms of connected Kan complexes over a general ring, homotopy fibers and Postnikov decomposition, and eventually equivalence of $Hot_n(\underline Z)$ with ordinary $Hot_n$. Page 465 distinguishes restriction of rings from extension, noting that restriction does not preserve projective components without an additional assumption and suggesting resolutions to remove that restriction. Its last sentence proposes compatibility of ring extension with the individual homotopy invariants; the next entry immediately revises that proposal. This transition is essential to reading the section as a developing argument rather than a list of established properties.

\subsection*{113. Soft Postnikov decomposition and a schematic fundamental group (pp. 466--472; 134-7/24--30).}
The entry dated 27.8 corrects the previous evening's assertion. Formula (A) concerns total homology and derived extension of scalars; the analogous termwise assertion (B) for homotopy modules fails without flatness hypotheses (p.~466). The revised arrow in (B) points left, a direction separately emphasized by a circled marginal arrow omitted by the arXiv witness. The reason is the use of shifted projective resolutions in a Postnikov fiber: after tensoring, higher homology can appear, so the base-changed successive fibration need not remain a Postnikov decomposition (p.~467). This explicitly answers the puzzle of p.~444 and motivates the judgment that ordinary Postnikov decomposition is not ``geometric,'' meaning compatible with extension of the ground ring. The proposed alternative keeps the entire abelianization complex. Using the completed divided-power linearization and the embedding $X\to W(L(X))$, Grothendieck proposes repeatedly taking the homotopy fiber of $X_*\to W(L(X_*))$, without first extracting the first nonzero homology module (pp.~467--468). He names this ``soft Postnikov d\'evissage.'' For a discrete $K(G,1)$ the analogy leads to iterated commutators, and he distinguishes the versions associated with solvability and nilpotence. Formula (D) contains an index problem: the visible inner exponent is $n+1$, like the outer exponent, while the arXiv witness gives $n-1$. The checked text preserves the scan's reinforced plus and does not silently substitute a mathematically expected recurrence. The next page adds the issue of the action of $\pi_1$ on higher homotopy groups; trivial or unipotent action appears relevant to unipotent schematization (p.~469). The second part of the section revisits complexes with $X_0=e$ but possibly nontrivial $X_1$. A group object gives a fundamental group that need not be abelian, and finite-presentation and smoothness assumptions lead to questions about nilpotence (p.~469). The Kan relations then motivate the presheaf $k'\mapsto\pi_1(X_*(k'))$. The desired invariant is a group object of $U(k)$ representing this sheaf, with an underlying parametrization by a module and a polynomial group law (pp.~470--471). This changes what the notation $\pi_1$ denotes: the group of sections alone is no longer the complete object. Formulas (G) and (G') express evaluation and base-change compatibility for the proposed schematic group. The following duplicated equation on p.~471 is nevertheless literally tautological in the scan, with $\pi_1(X_*(k))$ on both sides; the witness edits one side. The blank between $k$ and $k'$ in the preceding prose likewise contains no arrow on the page. To admit arbitrary quotient modules, the notion of submersion is reconsidered once more: maps isomorphic to those induced by epimorphisms of modules are proposed as the most convincing of the three possibilities considered so far (p.~471). The printed ``unless \ldots{} does'nt give rise'' contains an apparent extra negation, retained despite the witness's change to ``gives rise.'' The stronger requirement that the schematic fundamental group come from addition on a module is distinguished from commutativity; the discussion ends with tests involving the linearization theorem, a fibration over $K(G,1)$, and the relationship between nilpotent discrete groups and group objects of $U(k)$ (p.~472, referring back to section 94). The handwritten addition about unipotent action on the higher $\pi_i$ belongs here and is not merely incidental: it qualifies the contemplated range of homotopy types. Section 114 begins on the same page by describing the accumulated picture as provisional and preparing a review of the program.

\subsection*{114. A provisional model category and the completion gap (pp. 472--480; 134-7/30--38).}
This review arranges the preceding experiments into a prospective construction. The category $U(k)$ consists of unipotent bundles built from modules; smooth objects correspond to projective modules, submersions to epimorphisms, and smooth submersions to maps with projective kernel (pp.~472--473). The completed divided-power coalgebra describes the unpointed linearization, while the pointed version removes the completion in the basic construction (pp.~474--475). Completion is a substantive distinction: p.~474 explicitly says that the map (C) of p.~467 ``does'nt quite exist'' with $W$ as written, and introduces the completed version $W^{\wedge}$ instead. Base change for the completed functor consequently uses completed tensor products. The subsequent model categories impose $X_i=e$ for $i\leq n$, take sections to pointed semisimplicial sets, and localize at the proposed weak equivalences. The hoped-for comparison with ordinary homotopy types is restricted to the appropriate connected range over $\underline Z$ (p.~476). On p.~477 the first expression for the category of models visibly uses $\pi_n(k)$, although the surrounding formulas use the handwritten model-category symbol $\mathcal M_n(k)$; this notation anomaly is retained rather than silently regularized. Smooth Kan resolutions are then proposed to derive extension of scalars and the total linearization $LH.(X_*)=L_{pt}(X_*)$ (p.~477). The next two pages ask whether this homology functor is conservative and whether pointed linearization and completed linearization are quasi-isomorphic. They also reconsider comparison with ordinary free-module homology. The second question undergoes an immediate correction: on p.~479 the author says the proposed comparison with $k^{(X_*(k))}$ is generally false away from $k=\underline Z$. In the first nonvanishing degree, the expected $k$-module invariant is identified with integral homology of the section complex; tensoring that homology with $k$ produces a different candidate. The word \emph{integral} is explicit in the scan, whereas the arXiv witness gives \emph{internal}. The correction renews the need to construct homotopy modules intrinsically and to obtain them through a Postnikov argument, rather than transfer all ordinary homology assertions indiscriminately. The final page makes the outstanding construction precise. The exponential map gives $X_*\to W^{\wedge}(L(X_*))$, but the required homotopy-class map is $X_*\to W(L_{pt}(X_*))$ (p.~480). Even a quasi-isomorphism between the two linearizations does not by itself remove the difference between $W^{\wedge}$ and $W$; the handwritten phrase ``flaw in this `argument''' marks exactly this gap. The target of the completed map is a semisimplicial sheaf and need not belong to the proposed model category. Hard or soft Postnikov decomposition therefore still requires this factorization, a suitable notion of homotopy fiber, and homotopy constructions internal to the model category. The failure of $U(k)$ to have fiber products explains why Quillen's standard machinery cannot simply be applied as stated. These are explicit unresolved tasks, and section115 returns the following morning to the meaning of the coefficient sheaf itself.

\subsection*{115. Quasicoherent coefficients and prorepresentable linearization (pp. 481--489; 134-7/39--47).}
The entry dated28.8 restarts the formula numbering and identifies the source of the preceding confusion. The appropriate free linearization is over the quasicoherent ring sheaf $\underline O_k$ on the fpqc site, not over a constant ring (pp.~481--482). Grothendieck explains that the obstruction groups in schematic Postnikov decomposition use the quasicoherent sheaf associated with a homotopy module, rather than the same module regarded as a constant coefficient group. The circled marginal clarification of the argument $X(n)_*,\pi_{n+1}$ belongs to this distinction; the overwritten version in the running formula remains marked as doubtful. He connects the earlier mistake to having translated geometric language into commutative algebra too thoroughly, losing sight of which coefficient objects were intended. The universal free Module $\underline O_k^{(X)}$ gives the Hom, Ext and derived-Hom expressions (4)--(5') (p.~483), but is generally not quasicoherent. The proposed replacement would represent the same functor on quasicoherent coefficients. For $X=W(M)$, the continuous-Hom calculation through truncated divided-power modules shows that the correct object is a proobject, rather than the completed module regarded as an ordinary module (pp.~484--485). The obstruction is again scalar extension: inverse limits do not generally commute with tensoring by a new ground ring. Formula(13) packages the finite products of divided powers into a strict projective system; ordinary representability would require it to become essentially constant. The added phrase ``but with slightly different meaning'' therefore changes the interpretation of $L$ in (14), now valued in pro-modules. The pointed functor still takes values in ordinary modules, embedded as constant proobjects. Formula(18), on p.~486, supplies the exponential map into $WL(X)$ in the procategory and distinguishes this abelianization from its richer linear structure. The proposed derived comparison (20) and its Ext form are then tested on $X=W(M)$ (pp.~486--487). Formula(21) contains a meaningful handwritten revision from a limit expression to a direct sum of Ext groups of divided powers. Projectivity of $M$ gives the expected vanishing in positive degree, with a passage to finite-type submodules sketched for the infinite case. For nonprojective $M$, the potential nonzero Ext$^1$ term challenges the author's previous expectation that quasicoherent cohomology of these bundles should vanish. His reference to Larry Breen is a question about this expectation, not a proof of the general formula. The text visibly says ``quasiunipotent'' where the witness shortens it to ``unipotent'' (p.~487). The next page restricts attention to the smooth case and postpones the projective-resolution argument for full generality. For a semisimplicial complex, (22) uses Hyperext rather than a degreewise Ext construction; (23) compares the resulting pro-linearization with pointed linearization (p.~488). The comparison map(24) revisits the map from p.~454 with its proobject interpretation made explicit. Its relationship to a quasi-isomorphism, and then to the componentwise inverse limit used on p.~478, still requires foundational work. Page489 plainly admits that no argument has yet been found for the desired assertion under the proposed smoothness, connectivity and possibly Kan hypotheses. The second comparison is reformulated as $\underline O_k^{(X_*)}\to L(X_*)$ in (25). The author does not confidently assert that this is a quasi-isomorphism: quasicoherence of the first homology sheaves and essential constancy of the second remain in doubt. The practical hope is instead that it gives the needed comparisons against quasicoherent coefficients through(20). Section116 subsequently reports the collapse of the earlier general abelianization theorem, so these provisional claims must be read with that later revision in view.

\subsection*{116. The failed abelianization theorem and the interruption of the notes (pp. 490--491; 134-7/48--49).}
The entry dated 7.9 reports the decisive revision of the preceding program. After finding a simple description of schematic homotopy groups, Grothendieck learned in a telephone conversation with Illusie that the asserted comparison between discrete and schematic homology or cohomology was false. The anticipated equivalence between schematic and ordinary homotopy types therefore fails (p.~490). This is an explicit retraction of a central assumption underlying sections111--115; it does not cause him to abandon schematic homotopy types, whose formal properties he now wants to study in their own right, possibly through different models. The remainder of the section explains the intervening pause in writing through fatigue and the need to attend to the death of his granddaughter Ella. Page491 distinguishes work allowed to take place through attention from a compulsive forward movement that leaves perceptions disconnected from one's picture of reality. The revisions matter to the voice of this passage: the hand moves ``who'' after ``finally,'' retains ``his language'' for the body's intervention, and changes ``cause'' to ``caused.'' The witness substitutes ``which'' and ``its'' and loses the last correction. These pages provide both the mathematical warning needed to read the preceding conjectures and the author's own account of why the notebook paused before the technical review resumes in section117.

\subsection*{117. Size, exactness, and the distinction between schemes and formal schemes (pp. 492--496; 134-7/50--54).}
The entry dated 10.9 begins with two technical qualifications. The affine site used earlier is not automatically small in the working universe; a possible remedy is to bound the sizes of the modules and the rings used as arguments, using the fppf site for finite-type bundles (p.~492). More consequentially, $M\mapsto W(M)$ is right exact but need not preserve kernels. Thus a free resolution of $M$ need not become a resolution of $W(M)$. The derived-Hom comparison used to interpret formula(21) of p.~487 is therefore questioned in general and retained for the flat case (p.~493). This is a specific dependency of the cohomology calculation, not merely a set-theoretic caveat. The finite projective case clarifies the geometry through duality. Polynomial maps can be described by maps of symmetric algebras or by continuous maps of the dual completed coalgebras; preserving the origin translates into augmentation compatibility (pp.~493--494). The uncompleted divided-power coalgebra is instead dual to the adically completed symmetric algebra. Consequently, the two algebraic descriptions lead to schemes and formal completions along their zero sections, respectively. The scan's formulas require care: the prose on p.~493 includes a dual check on $M$, but formulas(2), (3) and (4) omit dual signs that the arXiv witness supplies. Those omissions remain in the checked text. The upward-arrow-shaped completion mark in(3') on p.~494 is likewise recorded as drawn, and a separate unreadable pen mark above the star in(5) remains marked. This geometric interpretation overturns the earlier hoped-for equivalence between the pointed and nonpointed linearizations (pp.~495--496). For the additive group in characteristic $p$, the degree-one comparison distinguishes additive polynomials from formal series $\sum_{i\geq0}c_i t^{p^i}$: the schematic side allows only finitely many nonzero coefficients. Grothendieck acknowledges that this example is neither simply connected nor over $\underline Z$, but states that he also no longer expects the comparison to be a quasi-isomorphism for the basic connected examples previously contemplated. He reports that Illusie rejected the vanishing assumptions behind his attempted ``way-out'' argument. The closing underline therefore warns explicitly against substituting $L_{pt}(W(M))$ for $L(W(M))$ when computing schematic homology or cohomology. Formal invariants have their own meaning and are distinct. Section118 begins on the same page with a different and deliberately striking claim: homotopy invariants themselves should nevertheless be recoverable from the first-order tangent object.

\subsection*{118. The tangent complex as a proposed homotopy invariant (pp. 496--500; 134-7/54--58).}
Against the distinction just established between schematic and formal cohomology, Grothendieck proposes that homotopy groups should depend only on the first-order tangent object. The Lie functor sends a smooth pointed bundle $W(M)$ to $M$, and formula(9) identifies the homotopy groups of $X_*(k)$ with those of the semisimplicial module $\mathrm{Lie}(X_*)$, hence with homology of its associated chain complex (p.~496). The next page sharply qualifies the claim: he has no direct description of a natural isomorphism under the stated hypotheses. It is tautological for complexes arising from chain complexes and is expected to extend through Postnikov decomposition, if every relevant model can be replaced by one admitting such a decomposition; naturality would still require work (p.~497). The typed revision makes this formula a ``cornerstone'' of a theory of schematic homotopy types and a proposed intrinsic definition of the invariants, rather than subordinating them from the outset to discrete homotopy groups. Weak equivalences would then be detected by the tangent complexes. Both subsequent comparison displays map an $X_*$ object to a $Y_*$ object; the arXiv witness repeats $X_*$ as target in each. The contemplated application is to explicit models of spheres (pp.~498--499). To make the Lie formula useful for ordinary homotopy groups, one must first realize the discrete type schematically over $\underline Z$ and obtain a workable model. Successive suspension from $S^1$ suggests the need for suspension constructions, complementing the earlier need for fibers and loops; the proposed derivator would organize the finite integration and cointegration operations. These are motivations and difficult construction problems, not completed calculations of sphere homotopy groups. The discussion of $S^2$ as a homogeneous space under $S^3$, and of the corresponding group/flag-variety geometry, asks for a model of the entire relationship: a group, a subgroup playing the role of a torus or Borel subgroup, and a quotient. Several overprinted group-name letters remain uncertain in the text, while the star on $\underline C^*$ is visibly an asterisk rather than the witness's multiplication cross. Hemispherical complexes supply a second route (pp.~499--500). Their smaller list of boundary and degeneracy operations, together with a prospective direct computation of homotopy in the discrete setting, is compared with the tangent-complex method. The marginal exclamation and vertical rule emphasize this alternative. The author immediately admits that he has not checked the proposed computation for arbitrary hemispherical sets that do not arise from an infinity-groupoid or infinity-stack. Comparing the two descriptions is therefore offered as a possible route toward understanding formula(9), with that unresolved premise preserved. Section119 then resumes on26.9 after another interruption and returns to the failure of the earlier homology comparison.

\subsection*{119. Integer-valued polynomials and a different comparison problem (pp. 500--505; 134-7/58--63).}
The entry dated 26.9 returns to the failed comparison of discrete and schematic linearization. Grothendieck reports Illusie's objection from Breen's work on stable cohomology, while admitting that the precise counterexample has not yet been made explicit to him (p.~500). He then gives a concrete degree-two obstruction using central extensions of the additive group (p.~501). The multiplicative group of truncated series $1+at+bt^2$ defines the extension $W_2$; a polynomial splitting requires $P(a+a')=P(a)+P(a')+aa'$. The solutions displayed over $\underline Q$ involve $t(t-1)/2$, which has no integral-coefficient representative after addition of the displayed linear term. This separates polynomial cochains over $\underline Z$ from arbitrary cochains on the discrete group. The final splitting display literally has $t$, without the square appearing in the preceding general homomorphism display; the checked text preserves that visible slip rather than adopting the arXiv witness's $t^2$. The main argument should therefore be read through the preceding coefficient equation, not by treating the last display as an independently correct formula. The next proposal replaces polynomial algebras by ``binomial algebras,'' built from the integer-valued expressions $P_n(t)=t(t-1)\cdots(t-n+1)/n!$ (p.~502). Grothendieck tentatively requires the ground ring to carry the corresponding binomial operations and compares this enrichment with divided powers in crystalline cohomology. The marginal reference to pp.~506--507 anticipates the next section's reconsideration of these structures and of the restrictions on the ground ring. On p.~503 the finite-type qualification on the free module $M$ is an above-line typed insertion. The algebra $\mathrm{Symbin}_{\underline Z}(M)$ is defined through rational polynomial functions integral-valued on the dual lattice, and the map to all functions on a semisimplicial free abelian group is posed as a possible quasi-isomorphism in the connected range. Reduction to $K(\underline Z,n)$ is proposed by d\'evissage, not proved. The broader question asks whether Eilenberg--Maclane homology or cohomology is derived from a suitable nonadditive functor $B$, leading to ``B-schematic homotopy types'' and unipotent $B$-bundles. The dark superscript on the first Symbin display remains unresolved and is not silently discarded. The author declines to pursue that alternative immediately, because the ordinary schematic program still promises a simple Lie-functor description of homotopy invariants (pp.~503--504). Loops, suspensions, and compatibility with passage to discrete types remain the practical tests. Reading Bousfield--Kan then changes the comparison: their construction is identified as $k$-completion, distinct from a fully developed category of types over $k$. The marginal ``definitely not, in general!'' rejects the tentative identification of completion with extension of the ground ring; the faint underline beneath ``ring extension'' identifies the disputed phrase (p.~504). Page~505 links this reading to the integration and cointegration operations of section~69 and to Don Anderson's existence theorem in closed model categories. A pen transposition moves the closed-model-category qualification directly after ``operations,'' making the setting of that theorem clearer. The section closes by recognizing Bousfield--Kan as foundational reading while postponing it long enough to record the recent schematic reflections. Section~120 starts on 28.9 by abstracting the extra algebraic operations just proposed.

\subsection*{120. Extra algebraic operations and polyhedral models (pp. 506--507; 134-7/64--65).}
The entry dated 28.9 abstracts the previous section's binomial proposal together with divided powers, lambda operations, and Adams-type operations. Grothendieck uses the notion of an ``analyzer'' and asks for an extension of the analyzer of commutative rings in which unary operations determine higher components by tensor powers. The handwritten changes from $n$ to $n+1$ tensor factors and from the older bounds to $n\geq-1$ and $n\geq1$ affect the indexing of the proposal (p.~506). The live corrected bounds are clear, although overwritten earlier digits remain doubtful. The essential datum is described as a ground ring $k_0$ and an algebra $\Omega_0$ equipped with composition. An $\Omega$-structure on a ring is then an action of these operations satisfying compatibility with addition, multiplication, and constants. Page~507 displays the two diagonals that encode the first two conditions. This is an outline of the desired axioms; the author expressly declines to write all of them down. The circled instructions ``\`a la ligne'' and ``id'' on p.~506 are retained as typist instructions, whose precise targets are not sufficiently explicit to justify speculative paragraph restructuring. The second half of p.~507 turns from these operations to David W.~Jones's work on polyhedral cells and poly-T-complexes, supplied by Ronnie Brown. Grothendieck relates it to the questions of sections~91 and~93: when does a category of cells form a weak test category, and when does the proposed standard chain complex compute the homology of its presheaves? He finds that Jones's emphasis is instead on generalizing thin structures and a corresponding equivalence of categories. The original questions therefore remain open for this notebook. Both appearances of the presheaf-category notation have a hand-drawn upward caret in the scan, rather than the downward vee supplied by the witness. Section~121 next returns to the formal role that linearization should play across the various possible schematic approaches.

\subsection*{121. Abelianization, adjunction, and the pointed correction (pp. 508--510; 134-7/66--68).}
Grothendieck begins the formal review by defining $\mathrm{Hotab}=D.(Ab)$ and considering abel\-ian\-iza\-tion $LH.$ together with the Eilenberg--Maclane functor $K_\pi$ (p.~508). He relates their adjunction to the free-abelian-group and forgetful functors on semisimplicial objects, then interprets ordinary cohomology through the representing objects $K(\pi,n)$. The notation is visibly under revision: a blacked-out earlier functor label is replaced by $K_\pi$, a handwritten $K_\pi$ is itself inserted and cancelled before the typed phrase ``defined via,'' and several category-name marks remain doubtful. The checked text does not reconstruct those older labels from the mathematical context. Page~509 explains why $K_\pi$ is conservative using $\pi_i(K_\pi(L.))\simeq H_i(L.)$, while explicitly denying full faithfulness and the assertion that the counit $LH.(K_\pi(L.))\to L.$ is an isomorphism. The typed letter $E$ in the preceding adjunction sentence is retained even though the surrounding notation uses $K_\pi$. The unit $X\to K_\pi(LH.(X))$ leads to the Hurewicz maps and a proposed homotopy-fiber interpretation of the Whitehead exact sequence (p.~509). The reduced homology in that sequence immediately calls for care with base points. Diagrams(8) and(10) factor the construction through pointed objects, using free pointing $X\amalg e$, the forgetful functor, and reduced free abelianization (pp.~509--510). The quotient description is replaced by an exact triangle, and the parenthetical typed ``one hopes!'' qualifies the continuing claim of adjunction. At the end of p.~510 the author explicitly corrects the earlier passage: the sequence and its fiber should come from the pointed unit(6'), since the unpointed expression lacks the canonical base points needed to define the indicated homotopy groups and fibers. The scan's(6') uses a tilde on both $K$ and $LH$, with no $\pi$ subscript on $K$; the witness regularizes this differently. Section~122 then revisits reduced abelianization itself, seeking a definition that no longer requires the chosen point.

\subsection*{122. Reduced abelianization without a point, and the failed homotopy-linearization functor (pp. 511--517; 134-7/69--75).}
The entry dated 2.10 revises the pointed construction of section~121. Instead of taking the quotient by the chosen base point, Grothendieck defines $\widetilde W(X)$ as the kernel of the augmentation $\underline Z^{(X)}\to\underline Z$ (p.~511). The fiber $W(X)(1)$ is the associated torsor and receives the universal map from $X$ to an abelian torsor. A point splits the extension, recovering the earlier reduced construction and giving the direct-sum description of homology; without a chosen point the natural exact triangle is arranged differently from the triangle on p.~510 (p.~512). The comparison is more informative over an object $Y$, because a section need not exist and the torsor class $c(X/Y)$ can be nontrivial. This supplies the connection with Postnikov invariants and motivates a more general Kan--Dold--Puppe correspondence for coefficients in the relevant slice topos (p.~513). The generalization to local test categories is announced as an intuition to revisit, not established here. The paper's notation is also simplified: the old $K_\pi$ and its pointed variant become $K$ and $\widetilde K$. A proposed partner $S$ or $\widetilde S$ would realize prescribed reduced homology, in contrast with the homotopy groups prescribed by $\widetilde K$ (pp.~514--515). The coefficient letter in the sphere-like objects $S(H,n)$ is changed in pen from $\pi$ to $H$, and circled marginal $H$ notes emphasize the change. Restriction to connected complexes addresses the immediate obstruction that degree-zero homology of a space is free. The author admits that he has neither constructed the desired functor nor established its existence or uniqueness. In the schematic setting even the relevant low-dimensional sphere objects have not yet been constructed. The possible value of this program depends on the tentative Lie-functor description from section~118, not on an available computation of sphere homotopy groups. The proposed fourth functor $L\pi.$ is more sharply rejected. It would turn homotopy groups into homology of a chain complex and satisfy $L\pi.(\widetilde K(L.))\simeq L.$. The continuation dated 3.10 states that these properties cannot all hold (p.~515). Applying the functor to the abelianization unit would provide a canonical section of the counit. On $L.=\pi[n]$, this would split the universal-coefficient sequence naturally in the groups involved (p.~516). The corrected indices matter: formula(18) has $\mathrm{Ext}^1$, $H^{n+1}$ and $H_{n+1}$, where the witness gives different indices. The text then specializes to degree one over $\underline F_2$ and the sequence $0\to V'\to\mathrm{Sym}^2(V')\to\bigwedge^2V'\to0$ (p.~517). An editorial mathematical qualification is needed at the last step. The manuscript claims equivariant nonsplitting for $\dim V\geq2$; the literal bound is retained. If $\mathrm{Aut}(V)$ means the finite group $\mathrm{Aut}_{\underline F_2}(V)$, as the preceding discussion of automorphisms of the abstract abelian group suggests, dimension two is an exception: $q=x^2+xy+y^2$ is fixed by all six invertible linear substitutions over $\underline F_2$, while its polar form maps to the generator of $\bigwedge^2V'$. It therefore supplies an equivariant section for that finite-group action. The qualification is specific to this interpretation: the calculation does not give equivariance for the algebraic group scheme $\mathrm{GL}(V)$ or for all matrices after scalar extension. This small calculation identifies a limitation of the asserted bound on the finite-group reading and does not repair the manuscript's text. The ensuing remark is conditional in a second way: it assumes that an analogous nonsplitting example can be found for $n\geq2$ before drawing the conclusion for simply connected homotopy types. Grothendieck presents that expected example as an alternative route to the failure of the earlier discrete/schematic comparison, avoiding reliance on the reported Breen results. Section~123 opens by acknowledging that, even after these discarded expectations, the basic schematic homotopy operations and the correct weak equivalences are still not in hand.

\subsection*{123. Sphere complexes and scalar operations as tests of the proposed theory (pp. 518--520; 134-7/76--78).}
The entry dated 4.10 takes stock after the preceding counterarguments. Even the equivalence of the proposed criteria for weak equivalence, and the construction of fibers and cofibers, are not yet settled (p.~518). Grothendieck therefore treats the realization of the ordinary two-sphere over $\underline Z$ as a basic test. Suspension should then produce the higher spheres; a larger category admitting the connected complexes needed for $S(1,k)$ might supply an alternative starting point (pp.~518--519). These are requirements for a workable theory, not constructions already obtained. If schematic sphere objects and the Lie description of their homotopy invariants exist, their homotopy groups should be the homology of canonical complexes $L\pi.(S^n_{\underline Z})$. Maps between spheres would then induce maps of those complexes, compatibly with composition (p.~519). The qualification from section~122 is crucial: the author has just rejected such a functor on arbitrary ordinary homotopy types. He now asks for a special structure on the sphere subcategory and explicitly allows the possibility that this structure is impossible, which would refute the intended schematic theory. The notation is still rough: the defining display for ${}^n\Pi.$ literally writes $L\,.(S^n)$, omitting the $\pi$ present in the surrounding discussion. The checked text retains that omission and the hand-corrected exponent and punctuation in the preceding displays. The first operation to examine is the action of the multiplicative monoid of integers, obtained by taking maps $S^n\to S^n$ (p.~520). The author relates the resulting actions on finite homotopy groups to the profinite monoid $\widehat{\underline Z}^{\mathrm{mult}}$ and proposes its lift to the derived base extension ${}^n\Pi.\otimes^{L}_{\underline Z}\widehat{\underline Z}$. The derived tensor symbol in the scan is not a direct-sum symbol, and no parentheses surround the tensor expression. For odd-dimensional spheres, restriction to the units is interpreted through the cyclotomic Galois action on the profinite homotopy type of punctured affine space. The pen changes from ``maximal'' to ``largest'' subgroup, adds ``(complex),'' and brackets the exponent $2m-1$; the above-line ``profinite'' remains typed. The point of the passage is the additional structure carried by these operations and a hoped-for lift to complexes, rather than a calculation of the sphere groups. Section~124 next separates the distinctive features of this schematic approach from the two alternatives considered earlier.

\subsection*{124. Four functors and a provisional formal calculus (pp. 521--530; 134-7/79--88).}
Grothendieck contrasts the present unipotent-bundle models with the earlier de Rham and diagonal-category approaches of sections~94 and~109. All seek a homology functor, but the present construction also paraphrases the Eilenberg--Maclane functor $\widetilde K$ compatibly with extension of scalars (p.~521). Simply forgetting a chain complex's $k$-module structure and taking a de Rham complex would not supply that compatibility. Moreover, the present schematic types map to ordinary types, whereas the earlier approaches had the reverse direction of comparison. The notation $\mathrm{Hotab}_0(k)$ deliberately leaves its algebraic realization open: section~115 required complexes of proobjects, and taking their inverse limits would lose the asserted base-change behavior. Consequently $\widetilde K$ may require pro-unipotent models on the nonlinear side as well. Pointed linearization offers a different formal theory, whose invariants section~117 has already distinguished from schematic ones (p.~522). The continuation dated 5.10 organizes four expected functors between $\mathrm{Hotab}_0(k)$ and $\mathrm{Hot}_0(k)$: total homology $L\widetilde H.$, total homotopy $L\pi.$, and their partners $\widetilde K$ and $\widetilde S$ (p.~523). Even the zero-connected convention is provisional: the author lacks an intrinsic satisfactory formulation of abelian fundamental group with trivial higher action, and briefly considers restricting to simply connected objects. Property~A posits the adjoint pairs $(L\widetilde H.,\widetilde K)$ and $(\widetilde S,L\pi.)$. Their units and counits give linearization and colinearization. The blank after $L$ in formula(5)b remains unfilled in the scan; the checked text does not silently add $\pi$. The Hom formulas interpret $K(\pi,n)$ as representing cohomology and $S(H,n)$ as representing mixed homotopy with a ``co-coefficient'' $H$. Taking $H=k$ recovers the expected ordinary homotopy-module functor (p.~524). This is the stated motivation for the spherical adjunction, whose existence has not been established. The circled marginal $H$ and $\widetilde W(X)$ record changes to the coefficient and linearization notation. Property~B gives the inversion identities $L\pi.(\widetilde K(L.))\simeq L.$ and $L\widetilde H.(\widetilde S(L.))\simeq L.$ (p.~525). Property~C obtains a Hurewicz map by either applying total homotopy to the linearization map or total homology to the colinearization map. The desired compatibility is that these constructions agree. The first nonvanishing homology objects should occur in the same dimension, with an isomorphism there; a pen addition also demands an epimorphism in the next dimension. The lengthy reduction to Eilenberg--Maclane homology in the preceding paragraph is present in the typescript even though the witness omits it. Property~D states preservation of cofibration sequences by the left adjoints and of fibration sequences by the right adjoints. The two exact triangles spell out the homology and homotopy cases (pp.~525--526). In the second triangle the scan's corrected top vertex is $L\Pi.(Y)$ and its bottom row is $L\Pi.(Z)\to L\Pi.(X)$; the witness gives a different cyclic arrangement. Both oblique labels are the numeral~1. The exactness discussion reveals an unwanted asymmetry. Suspension on the connected linear category is a shift and is fully faithful, whereas loops require truncation and lose information (p.~526). The first truncation bound has been corrected from~0 to~1. Remark~1 explains that the loop and fiber constructions therefore mean the neutral or marked connected component of the unrestricted construction. Grothendieck says that he has only just become aware of this restriction while writing, questions the practical value of imposing $X_0=e$, and contemplates enlarging both categories (p.~527). The wished-for trivial action of $\pi_1$ on higher groups is a different condition and is not supplied by $X_0=e$. Remark~2 recalls the Leray spectral sequence for homology of fibrations and asks whether a corresponding mechanism exists for homotopy of cofibers; it does not supply such a mechanism. Remark~3 starts from the evident exchange symmetry of the four-functor formulas, but then distinguishes autodual axioms from an actual autoduality of the categories (pp.~527--529). Full faithfulness of suspension must be excluded from the symmetric axioms. Derived duality on perfect complexes sends chains to cochains, prompting a tentative appeal to stabilization. The subsequent recollection of geometric duality is explicitly qualified by ``as far as I remember'' and ``if I got it right''; the former qualification, the word ``discovered,'' and the shift qualification are typed insertions, placed with pen brackets (p.~529). That discussion pertains to homology and still does not produce a duality exchanging homology and homotopy. The literal sentence ``homotopy type of X'' omits the vee used on the complement elsewhere; the checked text preserves the omission. Properties~E--G return to the operational requirements (pp.~529--530). Homology and homotopy should be conservative in the indicated connected setting, and all four functors together with their structural maps should commute with extension of the ground ring. Restriction of scalars is mentioned without a developed list of formal properties. Comparison with ordinary types combines forgetting $k$ on the linear side with taking sections of schematic models. It should respect $\widetilde K$, but the already disproved homology comparison is explicitly excluded, even for $k=\underline Z$. Compatibility with $\widetilde S$ is still hoped for over $\underline Z$, conditional on constructing the discrete spherical functor. There is no ordinary total-homotopy functor to compare with $L\pi.$; only the separate homotopy invariants are available there. The closing expectation that the comparison respects fiber and cofiber structures is another requirement of the program. Section~125 begins on13.10 with Baues's proposed axiomatic setting for making those operations precise.

\subsection*{125. From Baues's categories to weak equivalences and derivators (pp. 530--532; 134-7/88--90).}
The entry dated13.10 records a letter from H.~J.~Baues and a reading of his preprint on the homotopy classification problem. Grothendieck recognizes the relevance of the dual notions of cofibration and fibration category to his integration and cointegration operations (p.~530). The present semisimplicial unipotent models, however, do not admit the pushouts he would need for a cofibration structure, except in the stated trivial case. He has begun to consider a Kan-fibration notion but reports no definite conclusion. He therefore proposes to leave the schematic construction temporarily and return after further work on these categorical methods (p.~531). The more general objection concerns the scope of the axioms. Grothendieck questions whether every meaningful homotopy universe must arise from Quillen's or Baues's extra structures on a chosen category of models. He invokes semisimplicial objects in a topos, Illusie's weak equivalences, and Verdier's derived categories as reasons to seek a wider formulation (p.~531). These are the author's contemporary judgments and questions, not proofs of general nonexistence of model structures. The qualifications ``it seems,'' ``I would think,'' and ``but (it seems)'' are significant: the last is a pen revision to the claim about enough projectives or injectives and bounded derived categories (p.~532). The empty reference brackets after Illusie and the marginal cross remain empty in the text. The circled $C^{-}(A)$ and $D^{-}(A)$ correct the signs used in the discussion of the bounded categories. The proposed replacement starts with a category $M$ and weak equivalences $W$, without treating a selected fibration or cofibration class as the fundamental datum (p.~532). Different compatible choices should yield the same integration and cointegration operations. Section~69 supplies the earlier point of reference; a derivator is now identified as the structure that should express those operations independently of a particular model presentation. Constructing and characterizing the schematic derivator remain hopes. The passage also frames the derivator as the further structure missing when one stops at a localized or triangulated category, in both abelian and nonabelian settings. The typed insertions ``specific,'' ``on it,'' and ``make any mention of models'' sharpen that distinction, while the correction from ``is'' to the typed ``are'' refers to the two operations. Section~126 resumes the schematic discussion by testing the formal symmetry against Postnikov decomposition.

\subsection*{126. Postnikov truncation and the failure of a literal dual filtration (pp. 533--535; 134-7/91--93).}
The section tests the four-functor symmetry on the construction that originally motivated the schematic program. An $n$-co-connected object has no homotopy groups above degree $n$, and the proposed truncation $X\to X_n$ is the universal map to that full subcategory (p.~533). Equivalently, its inclusion should admit a left adjoint. This universal property produces the familiar inverse tower, with $X$ viewed as an appropriate inverse limit (p.~534). The discussion is axiomatic: it asks what should follow from the basic functors and their structural maps. The double underline on ``co'' distinguishes the upper-degree vanishing condition from connectedness. A literal dual would be a canonical increasing filtration ${}_nX\to X$, inducing the required homology isomorphisms through degree $n$ and vanishing above that degree. The ordinary skeletal filtration misses the endpoint of the desired isomorphism range. After trying to modify it, Grothendieck reports the feeling that no such canonical filtration exists, while explicitly admitting that he has neither formulated nor proved a nonexistence statement (p.~534). The decisive structural objection is that both proposed truncation subcategories are inverse images of the same linear subcategory under $L\pi.$ or $LH.$. That linear subcategory is stable under loops but not under suspension (p.~535). Even if the doubtful right adjoints existed on the homology side, their successive terms would not have the literally dual Postnikov relationship. A different suspension-stable family is contemplated and then abandoned for lack of a reasonable candidate. The conclusion qualifies the formal symmetry of section~124: homology and homotopy retain materially asymmetric properties. Several visible details prevent the reader from mistaking the formulas for a fully regularized exposition. On p.~533 the object named $Y$ is tested by a formula using $X$; on p.~534 the skeletal vanishing display writes $H_i(X_n)$ with a right subscript despite the surrounding left-subscript filtration. Formula labels are changed by hand, while some references to(5) remain on p.~535 after the relevant subcategory has been renumbered. Those literal features are retained. Section~127 then changes the question from symmetry of operations to the choice of components allowed in a schematic model.

\subsection*{127. Which objects should serve as schematic models? (pp. 535--538; 134-7/93--96).}
The section keeps the usual simplicial indexing category and changes the category of allowed components. The broad ambient setting is the fpqc topos of sheaves over the ground ring, with Illusie's weak equivalences of semisimplicial sheaves (pp.~535--536). Such a weak equivalence need not induce a weak equivalence on ordinary sections. Grothendieck therefore distinguishes the section functor from the canonical cointegration operation that should recover a discrete homotopy type. The above-line words ``canonical'' and ``(= weak equivalence)'' are typed revisions. He hopes that a suitably restricted class of models, especially the earlier unipotent bundles, will make the naive section construction adequate. The handwritten sentence opening p.~537 makes the base-change requirement explicit: a schematic homotopy type over $k$ should give a functor from $k$-algebras to ordinary homotopy types. The relative homotopy sheaves are defined first after choosing a section. The proposed ``pseudo-abelian'' condition combines relative zero-connectedness, abelian fundamental group, and trivial action on the higher homotopy sheaves (p.~537). Under a local formulation of that condition, the abelian sheaves should be independent of the chosen section and exist even without a global section. Their desired quasicoherent module structures are an additional issue. The comparison from the homotopy of sections to sections of the homotopy sheaves should be an isomorphism under mild hypotheses, for example $X_0=X_1=e$; the author expressly says he has not attempted a proof for unipotent bundles. He still calls the module structure a mystery without the Lie functor of section~118. The final remark proposes a spectral sequence from cohomology of the homotopy sheaves to the homotopy of cointegration, suggesting that acyclicity might imply the comparison (pp.~537--538). This remains a suggested argument, marked by ``maybe'' and ``If so,'' rather than a constructed spectral sequence. Section~128 next imposes a separate closure requirement: linearization should itself remain in the chosen category of bundles.


\chapter*{Sources}
Grothendieck, A. (1983). \emph{Pursuing Stacks} [Typescript with handwritten revisions; cotes 134-2--134-8]. Grothendieck archive, Université de Montpellier. The supplied page images and their metadata are the primary witnesses for this edition.

Grothendieck, A. (2021). \emph{Pursuing Stacks} (M. Carmona, with the collaboration of U. Buchholtz, Eds.; arXiv:2111.01000, version 2). Extended from the 2015 edition by ``the scrivener''; CC0 1.0. This edition is consulted as a comparison witness, not reproduced as the source of record.

Brown, R. (Trans.). (n.d.). \emph{English translation of the mathematical part of two letters from Alexander Grothendieck to L. Breen, 1975} [Typescript; prepared with assistance from L. Breen and J-L. Loday]. Grothendieck archive, Université de Montpellier, cote 134-2, sites 115--143.
\end{document}
